% ISEG MASTER DISSERTATION
% MASTER'S IN MATHEMATICAL FINANCE
% AUTHOR: Petr Terletskiy
% SUPERVISOR: [Advisor Name]
%
% Based on the ISEG template by Costa (2017),
% adapted by Domingos Guerreiro Seward and Luís F. Costa (2018).
%
% NOTE: Use XeLaTeX for compiling.
%
% =============================================
% Preamble
% =============================================

\documentclass[12pt,a4paper,oneside]{article}

% Font
\usepackage{times}
\usepackage[T1]{fontenc}
\usepackage[utf8]{inputenc}

% Language
\usepackage[english]{babel}

% Appendix
\usepackage[toc,page]{appendix}

\makeatletter
\def\@biblabel#1{\hspace*{-\labelsep}}
\makeatother

% Bibliography
\usepackage{natbib}

% Margins (3cm all sides per ISEG guidelines)
\usepackage[top=3cm, bottom=3cm, left=3cm, right=3cm]{geometry}

% Change margins command
\def\changemargin#1#2{\list{}{\rightmargin#2\leftmargin#1}\item[]}
\let\endchangemargin=\endlist

% Paragraph indentation
\usepackage{indentfirst}

% 1.5 line spacing
\usepackage{setspace}
\onehalfspacing

% Paragraph spacing
\parskip=6pt

% Headers and footers
\usepackage{fancyhdr}
\pagestyle{fancy}
\fancyhf{}
\lhead{\footnotesize \textsc{Petr Terletskiy}}
\rhead{\footnotesize\textsc{BTC Direction Prediction Using KANs}}
\cfoot{\thepage}

% Section formatting
\usepackage{titlesec}
\titleformat{\chapter}[display]
  {\scshape\center}
  {}
  {0pt}
  {\thechapter.\ }
\usepackage{etoolbox}
\makeatletter
\patchcmd{\chapter}{\if@openright\cleardoublepage\else\clearpage\fi}{}{}{}
\makeatother
\titleformat{\section}
  {\scshape\center}{\thesection}{1em}{}
\titleformat{\subsection}
  {\itshape\center}{\thesubsection}{1em}{}
\titleformat{\subsubsection}
  {\itshape\center}{\thesubsubsection}{1em}{}

% Hyperlinks
\usepackage{hyperref}
\hypersetup{
    colorlinks=true,
    linkcolor=black,
    citecolor=black,
    urlcolor=blue
}

% Table of contents
\usepackage{tocloft}
\setlength{\cftbeforetoctitleskip}{0pt}
\setlength{\cftbeforeloftitleskip}{0pt}
\setlength{\cftbeforelottitleskip}{0pt}

% Tables
\renewcommand*\thetable{\Roman{table}}
\addto\captionsenglish{\renewcommand{\tablename}{\textsc{Table}}}
\addto\captionsenglish{\renewcommand{\figurename}{\textsc{Figure}}}

\usepackage{lipsum}
\usepackage{booktabs}
\usepackage{adjustbox}
\setlength\heavyrulewidth{0.3ex}
\usepackage{tabularx}
\usepackage{array}
\usepackage{threeparttable}
\usepackage{multirow}
\usepackage{longtable}

% Quotes
\usepackage{csquotes}

% Mathematics
\usepackage{amsmath}
\usepackage{amsfonts}
\usepackage{amssymb}
\usepackage{bm}

% Graphics
\usepackage{graphicx}
\usepackage{float}
\usepackage{subcaption}

% Code listings (for algorithm pseudocode)
\usepackage{algorithm}
\usepackage{algpseudocode}

% Glossary / Acronyms
\usepackage[automake]{glossaries}
\makeglossaries

% ── Acronyms ──────────────────────────────────────────────────────────
% AFML & Methodology
\newacronym{afml}{AFML}{Advances in Financial Machine Learning}
\newacronym{tbl}{TBL}{Triple Barrier Labeling}
\newacronym{cusum}{CUSUM}{Cumulative Sum}
\newacronym{ffd}{FFD}{Fixed-width Window Fractional Differentiation}
\newacronym{cpcv}{CPCV}{Combinatorial Purged Cross-Validation}
\newacronym{dsr}{DSR}{Deflated Sharpe Ratio}
\newacronym{pbo}{PBO}{Probability of Backtest Overfitting}
\newacronym{mda}{MDA}{Mean Decrease Accuracy}

% Models
\newacronym{kan}{KAN}{Kolmogorov-Arnold Network}
\newacronym{lstm}{LSTM}{Long Short-Term Memory}
\newacronym{rf}{RF}{Random Forest}
\newacronym{xgb}{XGBoost}{Extreme Gradient Boosting}
\newacronym{lr}{LR}{Logistic Regression}
\newacronym{ar}{AR}{Autoregressive}
\newacronym{mlp}{MLP}{Multi-Layer Perceptron}

% Technical
\newacronym{btc}{BTC}{Bitcoin}
\newacronym{ohlcv}{OHLCV}{Open-High-Low-Close-Volume}
\newacronym{ta}{TA}{Technical Analysis}
\newacronym{ewma}{EWMA}{Exponentially Weighted Moving Average}
\newacronym{adf}{ADF}{Augmented Dickey-Fuller}
\newacronym{rsi}{RSI}{Relative Strength Index}
\newacronym{macd}{MACD}{Moving Average Convergence Divergence}
\newacronym{atr}{ATR}{Average True Range}
\newacronym{obv}{OBV}{On-Balance Volume}
\newacronym{mfi}{MFI}{Money Flow Index}
\newacronym{cci}{CCI}{Commodity Channel Index}
\newacronym{vix}{VIX}{CBOE Volatility Index}
\newacronym{mvrv}{MVRV}{Market Value to Realized Value}
\newacronym{sadf}{SADF}{Supremum Augmented Dickey-Fuller}

% Financial
\newacronym{mfw}{MFW}{Master's Final Work}
\newacronym{jel}{JEL}{Journal of Economic Literature}
\newacronym{emh}{EMH}{Efficient Market Hypothesis}


% =============================================
% Begin Document
% =============================================

\begin{document}

% ── Dedication ────────────────────────────────────────────────────────
\thispagestyle{empty}
\vspace*{\fill}
\begin{flushright}
\textit{%
% Add your dedication here, or leave blank.
}
\end{flushright}


% ── Glossary ──────────────────────────────────────────────────────────
\newpage
\thispagestyle{plain}
\addcontentsline{toc}{section}{Glossary}

\setcounter{page}{1}
\pagenumbering{roman}

\glsaddall
\printglossaries


% ── Abstract ──────────────────────────────────────────────────────────
\newpage
\thispagestyle{plain}
\addcontentsline{toc}{section}{Abstract, Keywords, and JEL Codes}

\section*{Abstract, Keywords, and \gls{jel} Codes}

% TODO: Write abstract (300 words max).
% State the problem, how you tackled it, and the main findings.
% Lead with results (Cochrane): say what you find, not what you look for.

\vspace{1em}

\textsc{Keywords}: Bitcoin; Kolmogorov-Arnold Networks; Symbolic Extraction; Triple Barrier Labeling; Combinatorial Purged Cross-Validation; Financial Machine Learning.

\textsc{\gls{jel} Codes}: C45; C53; C58; G11; G17.
% C45 - Neural Networks and Related Topics
% C53 - Forecasting and Prediction Methods
% C58 - Financial Econometrics
% G11 - Portfolio Choice; Investment Decisions
% G17 - Financial Forecasting and Simulation


% ── Table of Contents ─────────────────────────────────────────────────
\newpage
\thispagestyle{plain}
\addcontentsline{toc}{section}{Table of Contents}
\renewcommand*\contentsname{\hfill\normalfont\scshape\normalsize Table of Contents \hfill}
\tableofcontents

% ── List of Figures ───────────────────────────────────────────────────
\newpage
\thispagestyle{plain}
\addcontentsline{toc}{section}{List of Figures}
\renewcommand*\listfigurename{\hfill\normalfont\scshape\normalsize List of Figures \hfill}
\listoffigures

% ── List of Tables ────────────────────────────────────────────────────
\newpage
\thispagestyle{plain}
\addcontentsline{toc}{section}{List of Tables}
\renewcommand*\listtablename{\hfill\normalfont\scshape\normalsize List of Tables \hfill}
\listoftables

% ── Acknowledgements ──────────────────────────────────────────────────
\newpage
\thispagestyle{plain}
\addcontentsline{toc}{section}{Acknowledgements}
\section*{Acknowledgements}

% TODO: Write acknowledgements.


% =============================================
% TEXTUAL PART BEGINS
% =============================================
\newpage
\thispagestyle{plain}

\setcounter{page}{1}
\pagenumbering{arabic}

\begin{center}
\textsc{BTC DAILY DIRECTION PREDICTION USING KOLMOGOROV-ARNOLD NETWORKS WITHIN THE AFML FRAMEWORK}\\

\vspace{0.5cm}
By Petr Terletskiy\\

\begingroup\singlespacing
\begin{changemargin}{2cm}{2cm}
{\footnotesize \textsc{This dissertation} evaluates whether Kolmogorov-Arnold Networks can predict Bitcoin daily price direction and extract interpretable symbolic formulas, benchmarked against five models within López de Prado's Advances in Financial Machine Learning framework. Using Combinatorial Purged Cross-Validation with the Deflated Sharpe Ratio and Probability of Backtest Overfitting, the study finds that no model achieves statistically significant predictive performance, consistent with market efficiency. The core contribution is KAN symbolic expression extraction, deriving closed-form decision functions from the trained network.}
\end{changemargin}
\endgroup
\end{center}


% ==================================================================
% CHAPTER 1: INTRODUCTION
% ==================================================================
\section{Introduction}\label{sec:introduction}

% Target: max 20% of textual part (~7 pages for 35-page MFW).
%
% ── OPENING PARAGRAPHS (before any subsection) ───────────────────
% Cochrane: "Start with your central contribution."
% Thatcher R1: "First 1-2 sentences state a clear, field-relevant claim."
%
% Write 1-2 paragraphs that state the contribution upfront.
% Topics to cover:
%   - This thesis applies KANs to BTC daily direction prediction
%     within the full AFML framework
%   - Benchmarks KAN against five models (AR Logistic, LR, RF,
%     XGBoost, LSTM) under CPCV with DSR and PBO evaluation
%   - Core contribution: KAN symbolic expression extraction,
%     deriving closed-form decision functions from the trained network
%   - Telegraph the key finding: no model achieves statistically
%     significant performance (DSR=0), consistent with market efficiency;
%     the symbolic extraction pipeline successfully produces
%     interpretable formulas nonetheless


\subsection{Research Context and Contribution}\label{subsec:context}

% ── PROBLEM 1 (the spark): Interpretability and feature identification ──
%
% The narrative: "I wanted to understand what drives BTC price direction"
%
% Topics:
%   - ML models in finance are opaque; practitioners, risk managers,
%     and regulators cannot explain why a model made a specific prediction
%   - Researchers use dozens of features (23 TA indicators, 9 mathematical
%     features, 23 external: macro, crypto, on-chain) with no consensus
%     on which actually drive BTC direction
%   - On-chain data (active addresses, hash rate, MVRV, exchange flows)
%     offers a transparent window into blockchain network fundamentals
%     with no equivalent in traditional finance; the question is whether
%     this publicly available "inside view" translates into measurable signal
%   - The need: a model that not only predicts but also reveals WHAT
%     matters and HOW it matters, in a formula a human can read and audit
%
%
% ── PROBLEM 2 (the tool): Does model complexity justify its cost? ────
%
% The narrative: "My advisor pointed me to KANs as a solution"
%
% Topics:
%   - KANs (Liu et al., 2024) are a new neural network architecture based
%     on the Kolmogorov-Arnold representation theorem
%   - Unique property: learnable B-spline activation functions on edges
%     (not nodes), which can be distilled into closed-form symbolic
%     expressions via pruning and symbolification
%   - This makes KANs the only neural architecture that can produce
%     an explicit mathematical formula as its output, not just a prediction
%   - But KANs, LSTMs, and ensemble methods are orders of magnitude
%     more expensive than logistic regression on lagged returns
%   - If complex models don't produce better predictions under rigorous
%     evaluation, the interpretability contribution (symbolic formulas)
%     becomes the primary value, even under weak predictive performance
%   - No prior work applies KANs to BTC direction prediction, nor
%     extracts symbolic formulas from a classification KAN in a
%     weak-signal financial regime
%
%
% ── PROBLEM 3 (the realization): The evaluation crisis in crypto-ML ──
%
% The narrative: "When I surveyed the literature, I found it was
%                 methodologically broken"
%
% Topics:
%   - Most crypto-ML papers report 85-95% accuracy using naive
%     train/test splits with no purging, no embargo, fixed-horizon labels
%   - Overlapping triple-barrier labels leak information across folds,
%     inflating performance metrics; without purging and embargo,
%     reported results are unreliable
%   - The AFML framework (Lopez de Prado, 2018) was designed to prevent
%     exactly these pitfalls: CUSUM event filtering, TBL, FFD for
%     stationarity with memory preservation, sample uniqueness weights,
%     CPCV with purging and embargo, DSR, PBO
%   - Few crypto studies adopt the full AFML pipeline; most cherry-pick
%     individual components (e.g., TBL without CPCV, or CPCV without FFD)
%   - Applying the full framework honestly may produce negative results
%     (DSR=0, PBO near zero), which are not failures but evidence of
%     market efficiency under rigorous evaluation
%   - The three-version pipeline evolution in this thesis (early version
%     with data leakage, partial fix, final AFML-compliant version)
%     itself demonstrates how methodology shapes conclusions
%
%
% ── ASSET CHOICE: Why BTC? ───────────────────────────────────────────
%
% Topics (keep to one paragraph, factual justification only):
%   - BTC is the first and largest cryptocurrency by market capitalization,
%     representing over 50% of the total crypto market
%   - Since its inception in 2009, BTC has delivered the highest cumulative
%     returns of any asset class over the same period, including equities,
%     gold, and real estate
%   - It operates on a fully transparent, publicly auditable blockchain,
%     making it uniquely suited for on-chain feature analysis (this is
%     why on-chain features are part of the feature set)
%   - 24/7 trading, extreme volatility, and absence of market closes
%     create a demanding testbed for predictive models
%   - Daily OHLCV data freely available from 2014 onward provides
%     over a decade of history (~4,200 daily observations)
%
%
% ── CONTRIBUTION STATEMENT ───────────────────────────────────────────
%
% Topics (close the subsection with a crisp summary):
%   - This thesis fills three gaps by combining:
%     (a) The full AFML pipeline applied to BTC (evaluation rigor)
%     (b) KAN as a novel architecture for crypto direction prediction
%     (c) Symbolic expression extraction from a trained KAN
%   - Produces the first head-to-head comparison of six model
%     architectures under rigorous financial ML evaluation for BTC
%   - Produces the first symbolic formula extraction from a BTC
%     direction classifier
%   - Negative predictive results (DSR=0) are framed as honest
%     findings consistent with the efficient market hypothesis,
%     not as methodological failure


\subsection{Research Objectives}\label{subsec:objectives}

% Each RQ maps directly to a problem stated above.
% Thatcher R6: parallel logic (Problem -> RQ -> pipeline component -> result).
%
% Topics:
%
%   - RQ1 (from Problem 1 — interpretability + features):
%     Can a trained KAN produce interpretable symbolic decision
%     functions for BTC daily direction prediction, and which
%     features survive the extraction process?
%     → Answered by: symbolic extraction pipeline (Section 3.10),
%       results in Section 4.5
%
%   - RQ2 (from Problem 2 — complexity vs value):
%     Does KAN predictive performance justify its architectural
%     complexity relative to five established baselines (AR Logistic,
%     Logistic Regression, Random Forest, XGBoost, LSTM)?
%     → Answered by: CPCV model comparison (Section 4.1)
%
%   - RQ3 (from Problem 3 — evaluation crisis):
%     Does applying the full AFML methodology (CUSUM filtering,
%     triple barrier labeling, fractional differentiation, combinatorial
%     purged cross-validation, deflated Sharpe ratio, probability of
%     backtest overfitting) change the conclusions relative to the
%     naive evaluation approaches prevalent in the crypto-ML literature?
%     → Answered by: DSR and PBO results (Section 4.3),
%       three-version pipeline evolution (Section 5.1)


\subsection{Document Structure}\label{subsec:structure}

% One short paragraph. ISEG expects this; keep it minimal.
%
% Topics:
%   - Section 2 reviews the literature across three streams:
%     BTC price prediction, the AFML methodology, and KANs
%   - Section 3 presents the full methodology, from data collection
%     through labeling, feature engineering, CPCV, model training,
%     calibration, evaluation, and symbolic extraction
%   - Section 4 reports the results: model comparison, classification
%     and financial performance metrics, feature stability, and
%     the extracted symbolic formulas
%   - Section 5 discusses the findings, interpreting the negative
%     results in the context of market efficiency and positioning
%     the symbolic extraction contribution
%   - Section 6 concludes
%   - Section 7 outlines directions for future work


% ==================================================================
% CHAPTER 2: LITERATURE REVIEW
% ==================================================================
\newpage
\section{Literature Review}\label{sec:literature}

% Cochrane: "After you've explained your contribution, then you can write
% a brief literature review. Make it a separate section."
% Cochrane: "The main point should be to set your paper off against the
% 2 or 3 closest papers, and to give proper credit."
% Thatcher R3: Each subsection opens with a claim, not background.
%
% ── OPENING PARAGRAPH (before any subsection) ────────────────────
%
% Topics:
%   - One paragraph framing the three literature streams and how they
%     converge in this thesis
%   - This thesis sits at the intersection of three bodies of work:
%     (a) Kolmogorov-Arnold Networks as a novel interpretable architecture,
%     (b) ML-based cryptocurrency price prediction, and
%     (c) the AFML framework for rigorous financial ML evaluation
%   - Each stream is reviewed below, with emphasis on the gaps this
%     thesis addresses


% ------------------------------------------------------------------
\subsection{Kolmogorov-Arnold Networks}\label{subsec:lit_kan}
% ------------------------------------------------------------------

% Narrative: introduce the architecture, establish its unique capabilities,
% then show how it has been applied to finance, then identify the gap.
%
%
% ── 2.1.1 Mathematical Foundation ────────────────────────────────
%
% Topics:
%   - The Kolmogorov-Arnold representation theorem (Kolmogorov, 1957;
%     Arnold, 1958): any continuous multivariate function f(x1,...,xn)
%     can be decomposed as a finite sum of continuous univariate functions
%   - This means: any prediction task can theoretically be expressed as
%     compositions of one-variable functions, which is the foundation
%     for KAN's architecture
%   - The theorem was long considered impractical (the inner functions
%     can be highly non-smooth); Liu et al. (2024) made it practical
%     by parameterizing these functions with learnable B-splines
%
%
% ── 2.1.2 KAN Architecture (Liu et al., 2024) ───────────────────
%
% Topics (cite: liu_kan_2024):
%   - MLPs place fixed activation functions (ReLU, sigmoid) on nodes;
%     KANs place learnable activation functions on edges
%   - Each edge is parameterized as a B-spline: g(x) = sum(a_k * B_k(x))
%     where B_k are basis spline functions of order k over G grid intervals
%   - Architecture notation: [n_input, n_hidden, n_output]
%   - Grid extension: refine the spline grid from coarse to fine without
%     retraining from scratch (increases expressiveness)
%   - Training: typically Adam followed by L-BFGS for fine-tuning;
%     sparsity regularization (L1 on activations + entropy)
%   - Key advantage over MLPs: after training, each edge's spline
%     function can be inspected, pruned, or replaced with a symbolic
%     function (e.g., sin, exp, x^2), yielding an interpretable formula
%   - KANs achieve comparable accuracy to MLPs with significantly
%     fewer parameters on scientific discovery tasks
%
%
% ── 2.1.3 KAN 2.0 and MultKAN (Liu et al., 2024) ────────────────
%
% Topics (cite: liu_kan2_2024):
%   - KAN 2.0 introduces MultKAN: multiplication nodes alongside
%     addition, enabling multiplicative feature interactions
%     (e.g., rsi * stoch_k) without fragile log/exp decomposition
%   - Same symbolic extraction pipeline works for both standard
%     KAN and MultKAN
%   - Relevant to this thesis: MultKAN comparison is explored as
%     a potential extension (Section 7)
%
%
% ── 2.1.4 KANs for Financial Time Series ────────────────────────
%
% Topics:
%
%   VIX KAN paper (cite: cho_lee_kim_2025):
%   - Applies KAN to VIX forecasting (regression, not classification)
%   - Introduces Algorithm 1: Train (Adam + grid extend + LBFGS)
%     → Prune → Symbolify → Affine fine-tune
%   - Achieves competitive forecasting with interpretable closed-form
%     formulas; extracted formulas reveal mean reversion and leverage
%     effects in VIX, validating domain knowledge
%   - Direct methodological ancestor of this thesis's symbolic
%     extraction pipeline (Section 3.12)
%   - Key difference: VIX is a regression target on a relatively
%     predictable mean-reverting series; BTC daily direction is a
%     binary classification in a weak-signal regime, which makes
%     symbolic extraction substantially harder
%
%   KASPER (cite: oad_kasper_2025):
%   - Integrates KAN with regime detection (Gumbel-Softmax mechanism)
%     for stock price prediction
%   - Uses regime-specific KANs with sparse spline activations
%   - Achieves strong results (R^2 = 0.89, Sharpe = 12.02) on
%     individual stocks from Yahoo Finance
%   - Symbolic rule extraction via Monte Carlo Shapley values,
%     producing human-readable rules per regime
%   - Closest related work to this thesis in terms of KAN + finance
%     + interpretability
%   - Key differences from this thesis: (a) regression not classification,
%     (b) regime detection layer absent in our pipeline, (c) different
%     symbolic extraction method (Shapley-based vs direct formula),
%     (d) they do not use AFML evaluation methodology
%
%
% ── 2.1.5 Training Best Practices ───────────────────────────────
%
% Topics (cite: practitioner_kan_2025):
%   - Grid size selection: coarser grids generalize better on small
%     datasets; finer grids risk overfitting
%   - Regularization: L1 on spline activations + entropy regularization
%     encourages sparsity, which aids pruning and symbolification
%   - Adam vs L-BFGS staging: Adam for initial learning (escapes
%     local minima), L-BFGS for refinement (exploits curvature)
%   - Input normalization: tanh normalization to [-1, 1] matches the
%     default grid_range and stabilizes training
%   - These practices directly inform the training configuration
%     in this thesis (Section 3.6 and 3.12)
%
%
% ── 2.1.6 Broader Landscape ─────────────────────────────────────
%
% Topics (cite: kan_ts_review_2025):
%   - Brief mention of KAN variants for time series: TKAN (temporal),
%     ChebyKAN (Chebyshev basis), DecoKAN (decomposition), etc.
%   - The comprehensive review covers 40+ KAN-based models and notes
%     that symbolic distillation remains underexplored in financial
%     time series applications
%   - Cite as catch-all reference to avoid individually citing each
%     variant; this thesis uses standard B-spline KAN, not variants
%
%
% ── 2.1.7 Gap Statement ─────────────────────────────────────────
%
% Topics:
%   - No prior work applies KANs to BTC daily direction prediction
%   - No prior work extracts symbolic formulas from a classification
%     KAN (all existing work is regression: VIX forecasting, stock
%     price prediction)
%   - No prior work combines KAN symbolic extraction with the AFML
%     evaluation framework
%   - This thesis addresses all three gaps


% ------------------------------------------------------------------
\subsection{Bitcoin Price Prediction with Machine Learning}\label{subsec:lit_btc}
% ------------------------------------------------------------------

% Narrative: survey the BTC prediction landscape, show the range of
% approaches, then expose the methodological weaknesses that motivate
% adopting AFML in Section 2.3.
%
%
% ── 2.2.1 Overview of Approaches ────────────────────────────────
%
% Topics:
%   - BTC price prediction has attracted significant research interest
%     due to BTC's volatility, 24/7 availability, and growing
%     institutional adoption
%   - Three broad families of approaches:
%     (a) Statistical/econometric: ARIMA, GARCH, VAR
%     (b) Classical ML: Random Forest, XGBoost, SVM, Logistic Regression
%     (c) Deep learning: LSTM, GRU, CNN, Transformer, attention models
%   - Most studies frame the problem as either price-level regression
%     or direction classification (up/down); this thesis uses the latter
%
%
% ── 2.2.2 Technical Analysis and ML for BTC ─────────────────────
%
% Topics (cite: mate_confluence_2023):
%   - Investigates confluence of TA indicators with ML models for
%     BTC prediction
%   - Finds that ROC (Rate of Change) is among the top predictors,
%     which directly informs this thesis's feature set (roc_14)
%   - Demonstrates that combining multiple TA indicators outperforms
%     single-indicator strategies
%   - Limitation: uses standard train/test split without CPCV or
%     purging; does not account for label overlap
%
%
% ── 2.2.3 Deep Learning for BTC Direction ────────────────────────
%
% Topics (cite: dl_btc_direction_2024):
%   - Applies DL models (LSTM, CNN, hybrid architectures) to BTC
%     daily direction prediction, same problem as this thesis
%   - Reports competitive accuracy with LSTM-based architectures
%   - Does not use AFML labeling (fixed-horizon labels instead of TBL),
%     no CUSUM filtering, no sample weights
%   - No purging or embargo in cross-validation
%   - Key positioning: this paper solves the same problem with different
%     (and less rigorous) methodology; comparing conclusions illustrates
%     how evaluation methodology shapes reported performance
%
%
% ── 2.2.4 Broader Crypto-ML Landscape ───────────────────────────
%
% Topics (cite: crypto_dl_review_2024):
%   - Comprehensive review of DL models applied to cryptocurrency
%     forecasting across multiple coins and time horizons
%   - Documents the dominance of LSTM, GRU, and attention-based
%     models in the recent literature
%   - Notes recurring issues: small datasets, overfitting, lack of
%     out-of-sample robustness, inconsistent evaluation metrics
%   - Very few studies incorporate transaction costs, bet sizing,
%     or risk-adjusted performance metrics (Sharpe ratio, drawdown)
%   - Use as catch-all citation for the broader landscape; avoids
%     needing to cite 20+ individual crypto-ML papers
%
%
% ── 2.2.5 Methodological Weaknesses in the Literature ───────────
%
% Topics (this is the bridge to Section 2.3):
%   - Synthesize the common problems across the papers reviewed above:
%     (a) Naive train/test splits: chronological split without purging;
%         overlapping labels leak future information into training
%     (b) Fixed-horizon labeling: "did the price go up in exactly 5 days?"
%         ignores stop-loss and take-profit dynamics that real traders use
%     (c) No sample weighting: all observations treated equally, even
%         when labels overlap heavily (low uniqueness)
%     (d) Accuracy as the primary metric: ignores class imbalance,
%         ignores economic significance (a 52% accuracy model can
%         lose money after transaction costs)
%     (e) No correction for multiple testing: reporting the best model
%         from many trials without deflating the Sharpe ratio overstates
%         significance
%   - These are exactly the pitfalls that the AFML framework (next
%     subsection) was designed to address
%   - Position: this thesis applies AFML to produce a methodologically
%     honest evaluation of BTC direction prediction


% ------------------------------------------------------------------
\subsection{The AFML Methodology}\label{subsec:lit_afml}
% ------------------------------------------------------------------

% Narrative: present the AFML framework as the solution to the
% methodological problems identified in Section 2.2. Then show that
% few studies adopt the full pipeline, especially for crypto.
%
%
% ── 2.3.1 Overview and Motivation ───────────────────────────────
%
% Topics (cite: lopez_de_prado_2018):
%   - Lopez de Prado (2018) identifies systematic flaws in how ML
%     is applied to finance: backtest overfitting, information leakage
%     from overlapping labels, naive cross-validation, and inflated
%     Sharpe ratios from multiple testing
%   - Proposes a complete framework addressing each flaw:
%     data representation, labeling, sample weighting, feature
%     engineering, cross-validation, and backtesting evaluation
%   - The framework is designed for institutional equity trading but
%     the principles are asset-agnostic
%
%
% ── 2.3.2 Event-Driven Sampling: CUSUM Filter ──────────────────
%
% Topics (cite: lopez_de_prado_2018, Ch. 2):
%   - Standard time-sampled bars (daily OHLCV) produce observations
%     at fixed intervals regardless of market activity
%   - The symmetric CUSUM filter identifies structurally meaningful
%     events: timestamps where cumulative price deviation exceeds
%     a threshold h
%   - Reduces the dataset to informationally rich observations,
%     discarding periods of low activity
%   - Threshold calibration: h is typically set as a multiple of
%     mean daily volatility (this thesis uses h = 1.5x)
%
%
% ── 2.3.3 Triple Barrier Labeling ───────────────────────────────
%
% Topics (cite: lopez_de_prado_2018, Ch. 3):
%   - Fixed-horizon labeling ("did price go up in 5 days?") ignores
%     path-dependent risk management (stop-losses, take-profits)
%   - TBL defines three barriers for each event:
%     (a) Upper horizontal: take-profit at pt * daily_vol
%     (b) Lower horizontal: stop-loss at sl * daily_vol
%     (c) Vertical: maximum holding period (num_days)
%   - The label is determined by which barrier is touched first
%   - Produces labels that reflect real trading outcomes, not
%     arbitrary fixed-horizon snapshots
%   - Labels carry a natural t1 (barrier touch timestamp) needed
%     for downstream purging
%
%
% ── 2.3.4 Fractional Differentiation ────────────────────────────
%
% Topics (cite: lopez_de_prado_2018, Ch. 5):
%   - Integer differentiation (returns) achieves stationarity but
%     destroys long-term memory; price levels preserve memory but
%     are non-stationary
%   - FFD finds the minimum differentiation order d* that achieves
%     stationarity (ADF test) while preserving maximum memory
%   - Fixed-width window implementation truncates weights below
%     a threshold tau for computational efficiency
%   - d* is estimated per fold on training data only to prevent leakage
%
%
% ── 2.3.5 Sample Weights ───────────────────────────────────────
%
% Topics (cite: lopez_de_prado_2018, Ch. 4):
%   - Overlapping TBL labels mean some observations are more unique
%     (informative) than others
%   - Average uniqueness: for each label, compute 1/c_t averaged over
%     its lifespan, where c_t counts concurrent active labels at bar t
%   - Return-attribution weights: weight by |return| * uniqueness,
%     so high-return, unique events receive more influence
%   - Time decay: optionally downweight older observations to
%     prioritize recent market regimes
%   - These weights are passed to classifiers via sample_weight,
%     ensuring the model doesn't overfit to redundant observations
%
%
% ── 2.3.6 Combinatorial Purged Cross-Validation ─────────────────
%
% Topics (cite: lopez_de_prado_2018, Ch. 7 and 12):
%   - Standard k-fold CV is invalid for financial time series because
%     overlapping labels leak information across folds
%   - CPCV partitions data into N contiguous groups, selects k groups
%     as test set for each of C(N,k) combinations
%   - Purging: remove training observations whose labels overlap with
%     the test period (three overlap conditions)
%   - Embargo: remove training observations immediately after each
%     test group to prevent serial correlation leakage
%   - Path matrix: C(N-1,k-1) non-overlapping backtest paths allow
%     distribution-level analysis of performance (not just a single
%     point estimate)
%   - This thesis uses N=6, k=2: 15 splits, 5 backtest paths
%
%
% ── 2.3.7 Backtesting Evaluation: DSR and PBO ──────────────────
%
% Topics (cite: lopez_de_prado_2018, Ch. 11 and 14):
%
%   Deflated Sharpe Ratio:
%   - The observed Sharpe ratio is inflated when the best model is
%     selected from multiple trials
%   - DSR corrects for this selection bias by comparing the observed
%     Sharpe to the expected maximum under the null hypothesis
%     (all true Sharpes = 0)
%   - Accounts for non-normality (skewness, kurtosis) in returns
%   - DSR > 0.95 indicates the Sharpe is statistically significant
%     after correcting for the number of models tested
%
%   Probability of Backtest Overfitting:
%   - PBO measures how likely the in-sample best model underperforms
%     out-of-sample relative to the median model
%   - Computed via combinatorially symmetric cross-validation: split
%     the CPCV paths into IS and OOS halves, check if the IS winner
%     beats the OOS median
%   - PBO near 0 means in-sample selection is reliable; PBO > 0.5
%     means overfitting dominates
%
%
% ── 2.3.8 AFML Applied to Crypto ────────────────────────────────
%
% Topics (cite: slepaczuk_bieganowski_2024):
%   - One of the few studies combining FFD and TBL in crypto markets
%   - Applies supervised autoencoders with walk-forward FFD d*
%     estimation to BTC, ETH, and LTC
%   - Demonstrates that the AFML labeling + feature engineering
%     approach transfers to crypto with improved risk-adjusted returns
%   - Does not use CPCV (uses walk-forward instead), does not compute
%     DSR or PBO
%   - Key takeaway: validates FFD + TBL in crypto context, but
%     leaves CPCV evaluation as an open gap that this thesis fills
%
%
% ── 2.3.9 The Fitting Scheme Matters More Than Model Choice ─────
%
% Topics (cite: audrino_chassot_2024):
%   - Across 1,445 US equities, a properly fitted linear HAR model
%     outperforms RF, XGBoost, GBT, and FFNN for realized volatility
%     forecasting
%   - The key finding: studies that report ML superiority over HAR
%     used suboptimal fitting schemes (infrequent re-estimation,
%     short training windows) that handicapped the baseline
%   - When the fitting scheme is optimized (daily re-estimation,
%     2.5-4 year training window), the simple model wins
%   - Directly relevant to this thesis: our negative results (DSR=0)
%     are consistent with this finding; properly evaluated baselines
%     are hard to beat, especially in a weak-signal regime like BTC
%   - Also reinforces the importance of methodology over model
%     complexity, which is the thesis's Problem 3
%
%
% ── 2.3.10 Gap Statement ────────────────────────────────────────
%
% Topics:
%   - Few crypto-ML studies adopt the full AFML pipeline; most
%     cherry-pick individual components (e.g., TBL without CPCV,
%     or CPCV without FFD and sample weights)
%   - No study combines AFML + KAN for any asset
%   - No study reports DSR and PBO for KAN-based predictions
%   - This thesis is the first to apply the complete AFML pipeline
%     (CUSUM + TBL + FFD + sample weights + CPCV + DSR + PBO)
%     to BTC direction prediction with KANs and five benchmark models


% ==================================================================
% CHAPTER 3: METHODOLOGY
% ==================================================================
\newpage
\section{Methodology}\label{sec:methodology}

% Target: min 70% of textual part together with Results/Discussion.
% Present the pipeline sequentially, matching the code structure.
% Only present what the reader needs to understand the results.
% Cochrane: "There should be nothing before the main result that a
% reader does not need to know in order to understand the main result."
%
% Pipeline order: data → labeling → weights → features (+ alignment)
% → CPCV → preprocessing (FFD + scaling + selection) → models
% → training pipeline (calibration + tuning) → evaluation → symbolic


% ------------------------------------------------------------------
\subsection{Data}\label{subsec:data}
% ------------------------------------------------------------------
%
% Cochrane: "Do not do extensive data description of well-known datasets."
% BTC-USD daily OHLCV is well-known. Keep this section tight: source,
% range, quality controls, external sources overview. ~1 page max.
%
%
% ── Opening paragraph ─────────────────────────────────────────────
%
% Topics:
%   - The primary dataset consists of BTC-USD daily OHLCV prices
%     downloaded from Yahoo Finance via the yfinance Python library
%   - Date range: September 17, 2014, to April 12, 2026 (~4,200
%     daily observations)
%   - Start date chosen to begin after BTC's early illiquid period
%     (pre-2014 data is dominated by Mt. Gox and thin order books)
%   - Unlike equities, BTC trades 24/7 on global exchanges with no
%     weekend or holiday closures, so the daily bar calendar is
%     continuous
%
%
% ── Data quality ──────────────────────────────────────────────────
%
% Topics:
%   - Gaps: yfinance occasionally returns missing dates due to
%     exchange-level data issues; gaps of up to 3 calendar days are
%     forward-filled; gaps longer than 3 days raise an error (none
%     observed in the BTC-USD series)
%   - OHLCV consistency checks applied to every row:
%     (a) High >= max(Open, Close)
%     (b) Low <= min(Open, Close)
%     (c) Volume >= 0
%   - Rows with NaN Close are dropped before any downstream computation
%   - Index normalized to timezone-naive daily timestamps
%   - Adjusted prices used (auto_adjust=True in yfinance), so no
%     separate adjustment for splits is needed (BTC has no corporate
%     actions, but the flag ensures consistency)
%
%
% ── External data sources (brief overview) ────────────────────────
%
% Topics:
%   - Three categories of external data supplement the OHLCV features
%     (detailed in Section 3.4):
%     (a) Macro (13 features): traditional finance indicators from
%         yfinance and FRED (dollar index, treasury yields, VIX,
%         equity indices, commodities)
%     (b) Crypto-macro (2 features): cross-crypto market signals
%         from yfinance and CoinGecko (ETH/BTC ratio, BTC dominance)
%     (c) On-chain (8 features): blockchain network fundamentals
%         from CoinMetrics Community API (active addresses, hash rate,
%         MVRV, exchange flows, fees, supply)
%   - All external series are aligned to BTC's daily calendar using
%     merge_asof(direction='backward'), ensuring each BTC day uses
%     the most recent available value from each source without
%     look-ahead bias
%   - On-chain data is shifted by 1 day to account for CoinMetrics'
%     end-of-day reporting convention
%
%
% ── Figure suggestion ─────────────────────────────────────────────
%
%   Figure X: BTC-USD Close price (log scale) over the full sample
%   period (2014-2026), with CPCV group boundaries overlaid. This
%   figure can be reused from the notebook (cell 48) and serves
%   double duty: data overview + CPCV partition preview.


% ------------------------------------------------------------------
\subsection{Labeling}\label{subsec:labeling}
% ------------------------------------------------------------------
%
% This section describes the three-step process that converts raw
% price data into classified events: volatility estimation, event
% filtering, and label assignment. Each step is an AFML component
% with specific parameter choices that require justification.
%
% Thatcher R5 (thread the construct): use "event" consistently to
% mean a CUSUM-triggered timestamp, "label" for the TBL output,
% and "observation" for the final aligned sample.
%
%
% ── 3.2.1 Daily Volatility Estimation ────────────────────────────
%
% Topics:
%   - Exponentially weighted moving standard deviation of log returns
%     (AFML Snippet 3.1):
%       log_ret_t = log(Close_t / Close_{t-1})
%       daily_vol_t = EWMA_std(log_ret, span=50)
%   - The span parameter controls how quickly the volatility estimate
%     adapts to new information; smaller span = faster adaptation
%   - This thesis uses span=50 (vs Lopez de Prado's default of 100
%     for equities); justification:
%     (a) BTC exhibits faster regime transitions than equities
%         (halving cycles, regulatory shocks, sentiment-driven
%         rallies/crashes occur on shorter timescales)
%     (b) A 50-day window captures approximately 2 months of
%         activity, balancing responsiveness against noise
%     (c) Sensitivity analysis: span=50 produces a volatility
%         estimate that tracks BTC's well-known regime changes
%         (e.g., 2017 bubble, 2020 COVID crash, 2021 bull run)
%   - The daily volatility series has the same index as Close and
%     contains leading NaNs from the EWMA warm-up period
%   - This estimate serves two roles downstream:
%     (a) Target width for TBL horizontal barriers (Section 3.2.3)
%     (b) CUSUM threshold calibration (Section 3.2.2)
%
%   Cite: \cite{lopez_de_prado_2018}, Snippet 3.1
%
%
% ── 3.2.2 CUSUM Event Filter ─────────────────────────────────────
%
% Topics:
%   - Purpose: reduce the full daily timeline to a subset of
%     structurally meaningful events, discarding periods where
%     cumulative price change is insufficient to warrant a label
%   - Symmetric CUSUM (AFML Snippet 2.4):
%       s_pos = max(0, s_pos + log_ret_t)
%       s_neg = min(0, s_neg + log_ret_t)
%     An event fires when s_pos >= h or s_neg <= -h, then the
%     respective accumulator resets to zero
%   - Threshold: h = CUSUM_MULT * mean(daily_vol), where
%     CUSUM_MULT = 1.5
%   - Justification for 1.5x multiplier:
%     (a) A multiplier of 1.0 would fire too frequently (nearly
%         every other day), producing highly overlapping labels
%     (b) A multiplier of 2.0+ would produce too few events for
%         meaningful model training (< 500 events over 11 years)
%     (c) 1.5x is a moderate choice that balances event count
%         against informativeness; similar multipliers are used
%         in the AFML examples
%   - Effect on sample size: approximately 4,200 daily bars reduce
%     to approximately 900 CUSUM events (exact count depends on
%     the data range and should be reported from the actual run)
%   - The filtered event timestamps (t_events) serve as the index
%     for the TBL step below
%
%   Cite: \cite{lopez_de_prado_2018}, Snippet 2.4
%
%   Figure suggestion: CUSUM filter visualization showing the
%   positive and negative accumulators alongside BTC price, with
%   event markers at threshold crossings (notebook cell 15 has this)
%
%
% ── 3.2.3 Triple Barrier Labeling ────────────────────────────────
%
% Topics:
%   - For each CUSUM event at time t0, three barriers define the
%     label (AFML Snippets 3.2, 3.4, 3.5):
%
%     Upper barrier (take-profit):
%       p_upper = Close_{t0} * (1 + pt * daily_vol_{t0})
%     Lower barrier (stop-loss):
%       p_lower = Close_{t0} * (1 - sl * daily_vol_{t0})
%     Vertical barrier (time limit):
%       t1_max = t0 + num_days calendar days
%
%   - The label is determined by which barrier the price path
%     touches first:
%       +1 if upper barrier touched first (bullish event)
%       -1 if lower barrier touched first (bearish event)
%       sign(return) at vertical barrier if neither horizontal
%         barrier is touched, or 0 if |return| < min_return
%
%   - Parameter choices:
%     pt_sl = (1.5, 1.5): symmetric barriers
%       Justification: no directional bias; equal penalty for
%       missed upside and missed downside; symmetric barriers are
%       the standard in the AFML literature for unbiased evaluation
%
%     num_days = 10: approximately two trading weeks
%       Justification: short enough to capture actionable price moves
%       (a trader would not hold a daily-frequency position for
%       months), long enough to allow horizontal barriers to be
%       touched (with 1.5x vol barriers, a 5-day window would
%       rarely trigger; a 30-day window creates excessive label
%       overlap, complicating sample weights and purging)
%
%     min_return = 0.0: any non-zero return at the vertical barrier
%       receives a directional label
%
%   - The output DataFrame contains three columns:
%     ret: return at the first barrier touch
%     bin: class label {-1, 0, +1}
%     t1: timestamp when the label was resolved (first barrier touch)
%   - The t1 column is critical: it defines the label's temporal
%     span [t0, t1], which CPCV uses for purging (Section 3.5)
%
%   Cite: \cite{lopez_de_prado_2018}, Snippets 3.2, 3.4, 3.5
%
%   Figure suggestion: Triple-barrier visualization for 2-3 example
%   events showing the price path with the three barriers and the
%   touch point (notebook cell 18 has this)
%
%
% ── 3.2.4 Rare Label Removal ─────────────────────────────────────
%
% Topics:
%   - Classes representing fewer than 5% of total samples are
%     removed (AFML Snippet 3.8)
%   - In practice, with pt_sl=(1.5,1.5) and num_days=10, class 0
%     (vertical barrier touch with ambiguous direction) typically
%     falls below the 5% threshold and is dropped
%   - The result is a binary classification problem: {-1, +1},
%     mapped to {0, 1} inside the pipeline for sklearn compatibility
%   - Report the final class distribution (counts and proportions)
%     after rare label removal; approximate balance between Up and
%     Down is expected given the symmetric barrier configuration
%   - Final labeled event count: approximately 911 events (exact
%     count to be confirmed from the actual run)
%
%   Cite: \cite{lopez_de_prado_2018}, Snippet 3.8
%
%   Table suggestion: Label distribution table showing class counts,
%   proportions, and the number of events dropped as rare


% ------------------------------------------------------------------
\subsection{Sample Weights}\label{subsec:weights}
% ------------------------------------------------------------------
%
% Thatcher R3: open with the claim, not background.
% "Overlapping labels create redundancy that biases training."
%
% This section can be ~1 page. The four computation steps are
% sequential and can be presented as a single narrative flow.
%
%
% ── Opening: the problem ──────────────────────────────────────────
%
% Topics:
%   - TBL labels span from t0 to t1 (barrier touch), and these
%     intervals frequently overlap: multiple labels are "alive"
%     simultaneously at any given bar
%   - Training a classifier on equally-weighted overlapping labels
%     inflates the effective sample size: the model sees the same
%     market regime multiple times through different but correlated
%     labels
%   - This biases the model toward crowded periods (high overlap
%     = high effective weight) at the expense of unique, potentially
%     more informative events
%   - AFML Chapter 4 proposes a four-step weighting scheme that
%     corrects for this redundancy
%
%
% ── Step 1: Concurrent label count (Snippet 4.1) ─────────────────
%
% Topics:
%   - For each bar t in the full daily timeline, count the number
%     of active labels: c_t = |{i : t0_i <= t <= t1_i}|
%   - c_t is computed over the entire BTC daily index (not just
%     CUSUM events) to capture the true concurrency structure
%   - Bars with c_t = 0 fall outside any label's span
%   - Bars with high c_t indicate crowded periods where multiple
%     overlapping positions resolve simultaneously
%
%
% ── Step 2: Average uniqueness (Snippet 4.2) ──────────────────────
%
% Topics:
%   - For label i spanning [t0_i, t1_i]:
%       uniqueness_i = mean(1 / c_t) for t in [t0_i, t1_i]
%   - Interpretation: uniqueness near 1.0 means the label is the
%     only active label during its lifespan (highly informative);
%     uniqueness near 0 means many labels overlap (redundant)
%   - Labels that are entirely unique provide independent information
%     that the model should prioritize
%
%
% ── Step 3: Return-attribution weights (Snippet 4.10) ─────────────
%
% Topics:
%   - Combine the magnitude of the label's return with its uniqueness:
%       w_i = |ret_i| * uniqueness_i
%   - Normalization: weights are scaled so that sum(w) = len(w),
%     meaning the average weight equals 1.0; this is compatible
%     with sklearn's sample_weight convention
%   - Effect: high-return, unique events (rare, informative price
%     moves) receive the highest weights; low-return, overlapping
%     events (noise in crowded periods) receive the lowest
%   - This is the stronger variant from AFML; the simpler alternative
%     (uniqueness-only weights, Snippet 4.2) ignores return magnitude
%
%
% ── Step 4: Time decay (Snippet 4.11) ─────────────────────────────
%
% Topics:
%   - Linear decay applied chronologically: the oldest observation
%     receives a weight multiplied by oldest_weight = 0.5, the newest
%     receives a multiplier of 1.0, with linear interpolation between
%   - After decay, weights are re-normalized to maintain sum(w) = len(w)
%   - Justification: BTC's market microstructure evolves over time
%     (increasing institutional participation, new derivative products,
%     regulatory changes); weighting recent observations more heavily
%     prioritizes the current regime without discarding history entirely
%   - oldest_weight = 0.5: the oldest sample receives half the
%     influence of the newest; this is a moderate decay (1.0 = no decay,
%     0.0 = ignore old data entirely)
%
%
% ── Step 5: Outlier capping ───────────────────────────────────────
%
% Topics:
%   - After time decay, extreme weights are clipped at the 99th
%     percentile (weight_cap_quantile = 0.99)
%   - This prevents single extreme events (e.g., flash crashes with
%     very high |ret| and high uniqueness) from dominating training
%   - The number of capped weights is logged for transparency
%
%
% ── Integration with classifiers ──────────────────────────────────
%
% Topics:
%   - Sklearn models (LR, RF, XGBoost): weights passed via the
%     sample_weight parameter in .fit()
%   - Neural models (LSTM, KAN): weights integrated as per-sample
%     multipliers in the cross-entropy loss:
%       L = mean(w_i * CE(logits_i, y_i))
%     This is equivalent to sample_weight but implemented at the
%     loss function level for PyTorch compatibility
%   - Both approaches ensure that the model sees each observation
%     with influence proportional to its informativeness
%
%   Cite: \cite{lopez_de_prado_2018}, Chapter 4 (Snippets 4.1,
%   4.2, 4.10, 4.11)
%
%   Figure suggestion: Sample weights over time, with notable BTC
%   events annotated (COVID crash March 2020, 2021 ATH, 2022 FTX
%   collapse). Notebook cell 25 has this.


% ------------------------------------------------------------------
\subsection{Feature Engineering}\label{subsec:features}
% ------------------------------------------------------------------
%
% This section describes the 55 features organized in three groups:
% technical analysis (23), mathematical/microstructure (9), and
% external (23). A full feature list with descriptions belongs in
% Appendix E; the main text covers the design rationale, grouping
% logic, and anti-leakage measures.
%
% Thatcher R4 (one construct, one paragraph): each feature group
% gets its own block. Don't describe all 55 individually in prose;
% group them by function and reference the appendix table.
%
%
% ── Opening paragraph ─────────────────────────────────────────────
%
% Topics:
%   - 55 features organized in three groups reflecting different
%     information sources:
%     (a) Technical analysis (23): price and volume patterns
%         computed from OHLCV data (Section 3.4.1)
%     (b) Mathematical microstructure (9): AFML-inspired structural
%         features capturing entropy, complexity, stationarity,
%         and bubble dynamics (Section 3.4.2)
%     (c) External (23): macro, crypto-macro, and on-chain signals
%         from sources outside the BTC price series (Section 3.4.3)
%   - All features are computed over the full daily timeline, not
%     just at CUSUM events; restriction to labeled events happens
%     later in the alignment step (Section 3.4.5)
%   - All features are backward-looking: each value at time t depends
%     only on information available at or before t
%
%   Table suggestion: Summary table with three rows (TA, Math,
%   External), columns for count, source, and purpose. Detailed
%   per-feature table in Appendix E.
%
%
% ── 3.4.1 Technical Analysis Features (23) ────────────────────────
%
% Topics:
%   - Computed from OHLCV using standard TA formulas; no external
%     data required
%   - Grouped by function (not alphabetically) for readability:
%
%     Returns and volatility (6):
%       log_returns: log(Close_t / Close_{t-1}); the fundamental
%         return measure; also used by AR Logistic as input
%       realized_vol: rolling 21-day std of log_returns, annualized
%         (* sqrt(365)); captures short-term volatility level
%       gk_vol: Garman-Klass (1980) volatility estimator; uses
%         OHLC data for more efficient estimation than close-to-close
%       yz_vol: Yang-Zhang (2000) volatility; best unbiased OHLC
%         estimator, combines overnight, open-to-close, and
%         Rogers-Satchell components
%       atr: Average True Range (EWMA smoothed, span=14); measures
%         typical price range per bar; log-transformed (Section 3.4.4)
%       bb_width: Bollinger Band width (20-day); dimensionless
%         measure of price dispersion relative to its moving average
%
%     Momentum and trend (9):
%       rsi: Relative Strength Index (14-day, Wilder smoothing)
%       macd, macd_signal, macd_hist: MACD system (12/26/9)
%       roc_14: Rate of Change; identified as a top predictor for
%         BTC in the Confluence paper (Mate et al., 2023)
%       stoch_k, stoch_d: Stochastic oscillator (14-day, 3-day smooth)
%       williams_r: Williams %R (momentum confirmation)
%       cci_14: Commodity Channel Index (14-day)
%
%     Volume (3):
%       obv: On-Balance Volume; cumulative, log-transformed
%         (Section 3.4.4) to compress the exponential growth
%       chaikin_osc: Chaikin Oscillator (3/10 EMA of ADL);
%         volume-weighted price momentum
%       mfi_14: Money Flow Index (14-day); volume-weighted RSI
%
%     Distribution shape (2):
%       skewness: rolling 21-day skewness of log returns
%       kurtosis: rolling 21-day excess kurtosis of log returns
%
%     Trend ratios (3):
%       ema_ratio_20_50: short vs medium trend strength
%       ema_ratio_50_200: golden/death cross signal
%       vwma_ratio_20_50: volume-weighted trend confirmation
%
%   - All rolling windows use 21 days (~1 trading month for BTC)
%     unless the indicator has a standard period (RSI=14, MACD=12/26/9)
%   - Parameters follow standard TA conventions; no optimization
%     of indicator periods is performed (this is deliberate: optimizing
%     TA parameters would introduce another layer of overfitting risk)
%
%   Cite: \cite{garman_klass_1980}, \cite{yang_zhang_2000},
%   \cite{mate_confluence_2023} for ROC as top BTC predictor
%
%
% ── 3.4.2 Mathematical Features (9) ──────────────────────────────
%
% Topics:
%   - AFML Part 4 introduces structural features that capture
%     properties of the return-generating process itself, rather
%     than specific price or volume patterns
%   - These features are computationally expensive (O(n^2) for
%     SADF and SMT) and cached to Parquet after first computation
%
%     Information-theoretic (3):
%       shannon_entropy (window=21): quantile-encoded Shannon entropy
%         of log returns; measures randomness/predictability; higher
%         values = more random, lower = more structured
%       lz_complexity (window=63): Lempel-Ziv-76 complexity;
%         measures compressibility of the binary return sequence
%         (positive/negative); normalized by window/log2(window)
%       gaussian_entropy (window=21): H_gauss = 0.5*ln(2*pi*e*var);
%         the entropy a Gaussian with the same variance would have;
%         the gap between this and Shannon entropy quantifies
%         non-Gaussianity of the return distribution
%
%     Random walk tests (2):
%       hurst (window=126): R/S Hurst exponent estimated via
%         rescaled range analysis at sub-periods [10, 21, 42, 63];
%         H > 0.5 indicates trending/momentum, H < 0.5 indicates
%         mean reversion, H = 0.5 is consistent with a random walk
%       variance_ratio (window=63, lag=5): Lo-MacKinlay (1988)
%         VR(q) = Var(r_q) / (q * Var(r_1)); VR > 1 indicates
%         momentum, VR < 1 indicates mean reversion, VR = 1 is
%         consistent with a random walk
%
%     Normality test (1):
%       jarque_bera (window=63): JB = (n/6)*(S^2 + K^2/4);
%         measures departure from normality; higher values indicate
%         stronger fat tails and asymmetry (tail risk indicator)
%
%     Structural break detection (3):
%       sadf: Supremum ADF (AFML Snippet 17.1); the maximum ADF
%         t-statistic over all expanding sub-windows starting from
%         SADF_MIN_SL=63; detects explosive behaviour (bubbles)
%         in log prices
%       smt_poly1, smt_exp: Sub/Super-Martingale tests (AFML
%         Section 17.4.3); supremum of normalized t-statistics for
%         polynomial and exponential trend specifications; detect
%         whether the price process departs from a martingale
%
%   Cite: \cite{lopez_de_prado_2018} Part 4, \cite{lo_mackinlay_1988}
%
%
% ── 3.4.3 External Features (23) ─────────────────────────────────
%
% Topics:
%   - Three sub-groups representing different information layers
%     beyond BTC's own price and volume history:
%
%     Macro (13): traditional finance signals reflecting the
%     broader economic environment
%       dxy_roc_21: US Dollar Index 21-day rate of change (%);
%         dollar strength inversely correlated with BTC in many periods
%       us2y, us10y: Treasury yields (levels); from FRED DGS2
%         (with T10Y2Y spread fallback if FRED is unavailable)
%         and yfinance ^TNX
%       yield_curve_2y10y: us10y - us2y; slope of the yield curve;
%         inversion signals recession expectations
%       yield_curve_10y30y: us30y - us10y; long-end term premium
%       vix: CBOE Volatility Index (level); "fear gauge" for
%         equity markets, often inversely correlated with risk assets
%       sp500_ret_21, nasdaq_ret_21: 21-day log returns of S\&P 500
%         and Nasdaq; capture equity market momentum
%       gold_ret_21, silver_ret_21, copper_ret_21: commodity
%         momentum; gold as safe haven, copper as growth proxy
%       oil_ret_21, natgas_ret_21: energy momentum; energy costs
%         affect mining economics and broader risk appetite
%
%     Crypto-macro (2): market-level cross-crypto signals (not
%     blockchain fundamentals)
%       eth_btc_ratio: ETH price / BTC price; captures relative
%         strength between the two largest cryptocurrencies;
%         falling ratio suggests BTC dominance
%       btc_dominance: BTC market capitalization share; fetched
%         from CoinGecko (with ETH/BTC inverse proxy as fallback);
%         the absolute market cap carries the same signal after
%         feature scaling
%
%     On-chain (8): blockchain network activity from CoinMetrics
%     Community API
%       active_addr_roc_14: 14-day rate of change in active
%         addresses; rising = growing network adoption
%       tx_count_roc_14: 14-day rate of change in transaction count;
%         rising = increasing usage
%       hashrate_roc_30: 30-day rate of change in hash rate;
%         rising = increasing mining investment / network security
%       mvrv: Market Value to Realized Value ratio (level);
%         MVRV > 1 = market overvalued relative to cost basis,
%         MVRV < 1 = undervalued; mean-reverting, already stationary
%       net_exchange_flow: daily exchange inflows minus outflows (BTC);
%         positive = coins moving to exchanges (selling pressure),
%         negative = coins leaving exchanges (accumulation)
%       fee_per_tx: total fees / transaction count; proxy for
%         network congestion and user willingness to pay for
%         block space
%       exchange_supply_pct: (supply on exchanges / total supply) * 100;
%         declining trend suggests long-term holding behaviour
%       issuance_ntv: daily BTC issuance; changes discretely at
%         halving events (approximately every 4 years)
%
%   - Anti-leakage measures for external data:
%     (a) All external series aligned to BTC's daily calendar via
%         merge_asof(direction='backward'): each BTC day uses the
%         most recent available value from the external source,
%         never a future value
%     (b) CoinMetrics data shifted by 1 day inside the fetch function:
%         CoinMetrics reports end-of-day aggregates, so the raw
%         value at date t summarizes activity through t; shifting
%         ensures the model sees this only from t+1 onward
%     (c) FRED treasury data falls on business days only; merge_asof
%         forward-fills the last available business day value into
%         weekends and holidays
%
%   Cite: CoinMetrics Community API (free tier, no API key required)
%
%
% ── 3.4.4 Log Transforms ─────────────────────────────────────────
%
% Topics:
%   - Two features require log transformation to compress their
%     scale before downstream scaling:
%     (a) ATR: log(|x| + 1e-8); ATR is always positive and varies
%         by orders of magnitude across BTC's price history
%         ($100 ATR in 2015 vs $5,000 ATR in 2024)
%     (b) OBV: sign(x) * log(|x| + 1); OBV is cumulative and grows
%         exponentially; the signed log preserves the direction
%         of volume flow while compressing the scale
%   - Log transforms are applied after feature computation and
%     before alignment; they affect only these two columns
%
%
% ── 3.4.5 Alignment ──────────────────────────────────────────────
%
% Topics:
%   - The feature matrix (55 columns, ~4,200 rows covering every
%     daily bar) must be restricted to the ~911 CUSUM-filtered events
%     that have labels and weights
%   - Alignment computes the date intersection of three objects:
%     features (full daily index), labels/bins (CUSUM subset), and
%     sample weights (same CUSUM subset)
%   - Produces four aligned objects that enter the CPCV loop:
%     X: pd.DataFrame (911 x 55) feature matrix
%     y: pd.Series (911) class labels {-1, +1}
%     w: pd.Series (911) sample weights
%     t1: pd.Series (911) barrier touch timestamps
%   - Post-alignment validation checks:
%     (a) No duplicate dates
%     (b) Monotonically increasing index
%     (c) No entirely NaN feature columns
%     (d) t1 completeness (no NaN values; required for CPCV purging)
%     (e) Shape consistency across X, y, w, t1
%   - Remaining NaN values in individual feature columns (from
%     rolling window warm-up or external data gaps) are logged
%     as warnings; they are handled by forward/back-fill inside
%     the per-fold preprocessing step (Section 3.6)
%   - The aligned sample size (~911 observations) is the effective
%     dataset for all model training and evaluation
%
%   Table suggestion: Alignment summary showing total daily bars,
%   CUSUM events, labeled events, aligned sample size, date range,
%   class distribution


% ------------------------------------------------------------------
\subsection{Cross-Validation Framework}\label{subsec:cpcv}
% ------------------------------------------------------------------
%
% Thatcher R3: open with the claim.
% "Standard k-fold CV is invalid for financial time series with
% overlapping labels; CPCV prevents the resulting information leakage."
%
% This section is critical: CPCV is the backbone of the evaluation.
% A committee member who doesn't understand purging and embargo
% cannot evaluate the results. Invest 1.5-2 pages here.
%
%
% ── Opening: why standard CV fails ────────────────────────────────
%
% Topics:
%   - In standard k-fold CV, observations are randomly assigned to
%     folds; the model trains on fold A and predicts fold B
%   - This is invalid when labels overlap in time: a TBL label at
%     t0 with barrier touch at t1 spans a multi-day interval [t0, t1]
%   - If a training observation's label span overlaps with a test
%     observation's label span, the training set contains information
%     about the test period's price path
%   - This leakage inflates metrics: the model appears to predict
%     the future, but it is merely recognizing patterns from
%     overlapping training labels
%   - CPCV (Lopez de Prado, 2018, Chapters 7 and 12) addresses this
%     through three mechanisms: combinatorial splitting, purging,
%     and embargo
%
%   Cite: \cite{lopez_de_prado_2018}, Chapters 7 and 12
%
%
% ── CPCV configuration ───────────────────────────────────────────
%
% Topics:
%   - The aligned dataset of T observations (~911) is partitioned
%     into N = 6 contiguous chronological groups
%   - Group construction: groups 0 through 4 contain floor(T/6)
%     observations each; group 5 absorbs the remainder
%   - For each split, k = 2 groups are designated as the test set,
%     and the remaining N - k = 4 groups form the initial training set
%   - Total splits: C(N, k) = C(6, 2) = 15
%   - Each group appears in C(N-1, k-1) = C(5, 1) = 5 test sets
%   - The 15 splits generate phi[N, k] = C(N-1, k-1) = 5 non-overlapping
%     backtest paths; each path covers all N groups exactly once,
%     drawing predictions from different splits
%
%   - Justification for N=6, k=2:
%     (a) N=6 with ~911 observations yields ~152 observations per group,
%         providing enough test data per fold for meaningful metrics
%     (b) k=2 keeps the test set at ~304 observations (~33% of data),
%         leaving ~607 for training (~67%); this ratio balances
%         training capacity against test set statistical power
%     (c) 15 splits and 5 paths provide sufficient combinatorial
%         diversity for PBO computation without excessive computation
%     (d) Increasing N would shrink groups and reduce per-group
%         sample size; decreasing N would reduce the number of
%         splits and paths available for distributional analysis
%
%   Table suggestion: Group boundaries table showing group ID,
%   positional index range [start, end), date range, and number
%   of observations per group
%
%
% ── Purging (AFML Snippet 7.1) ───────────────────────────────────
%
% Topics:
%   - For each of the 15 splits, training observations whose labels
%     temporally overlap with the test period are removed
%   - Three overlap conditions are checked for each training
%     observation i with label span [t0_i, t1_i] against each
%     test group spanning [t_test_start, t_test_end]:
%
%     Condition 1: training observation starts within test period
%       t_test_start <= t0_i <= t_test_end
%       (the observation itself falls in the test window)
%
%     Condition 2: training label resolves within test period
%       t_test_start <= t1_i <= t_test_end
%       (the label's barrier touch timestamp falls in the test
%       window, meaning the price path during the test period
%       determined this training label's outcome)
%
%     Condition 3: training label spans the entire test period
%       t0_i <= t_test_start AND t_test_end <= t1_i
%       (the label's span encompasses the full test window)
%
%   - Any training observation satisfying at least one condition
%     is purged from the training set for that split
%   - Purging is the primary defense against information leakage;
%     without it, the model could learn from training labels whose
%     outcomes depend on test-period price movements
%
%   Cite: \cite{lopez_de_prado_2018}, Snippet 7.1
%
%
% ── Embargo (AFML 7.4.2) ────────────────────────────────────────
%
% Topics:
%   - After purging, a buffer zone of training observations
%     immediately following each test group is removed
%   - embargo_pct = 0.01 (1% of T); with T ~= 911, this removes
%     approximately 9 observations after each test group boundary
%   - Purpose: prevent serial correlation leakage; even after
%     purging, training observations immediately after the test
%     period may be correlated with test observations through
%     shared market microstructure (autocorrelation in volatility,
%     momentum spillover)
%   - Embargo is applied only after the test group (not before),
%     following AFML 7.4.2: training labels that resolve before
%     the test period begins contain no future information about
%     the test period and should not be removed
%   - The combined effect of purging + embargo reduces the training
%     set by approximately 10-30 observations per split (depending
%     on label overlap density near group boundaries)
%
%   Cite: \cite{lopez_de_prado_2018}, Section 7.4.2
%
%
% ── Path matrix (AFML 12.4.1) ────────────────────────────────────
%
% Topics:
%   - The 15 splits and 6 groups produce 5 backtest paths
%   - Each path assigns each group to exactly one split, such that
%     the path's predictions cover the entire timeline without gaps
%     or overlaps
%   - Construction: for each group g, there are exactly 5 splits
%     that include g in their test set; path p uses the p-th such
%     split for group g
%   - The result is 5 independent backtest trajectories, each
%     producing its own Sharpe ratio, drawdown, and cumulative return
%   - This enables distribution-level performance analysis: instead
%     of reporting a single Sharpe ratio (which is a point estimate
%     with unknown sampling error), the thesis reports the median,
%     standard deviation, and range of Sharpe ratios across 5 paths
%   - The path Sharpe matrix (n_models x n_paths) feeds into both
%     DSR and PBO computation (Section 3.9)
%
%   Cite: \cite{lopez_de_prado_2018}, Section 12.4.1
%
%   Figure suggestion: CPCV split visualization showing 3
%   representative splits (e.g., splits 0, 6, 14) with train/test
%   regions, purged zones, and embargo zones marked. Notebook cell 50.
%
%
% ── Leakage verification ─────────────────────────────────────────
%
% Topics:
%   - An empirical leakage audit verifies that purging and embargo
%     work correctly across all 15 splits
%   - For each split, the audit checks: are there any training
%     observations whose t1 (barrier touch timestamp) falls within
%     a test group's date range?
%   - Result: zero leaks detected across all 15 splits
%   - This is reported as a table (split ID, test groups, train size,
%     test size, leak count, status) to provide transparent evidence
%     that the anti-leakage machinery operates as designed
%   - The audit is a deliberate methodological safeguard: it verifies
%     the implementation empirically, not just the algorithm's
%     theoretical correctness
%
%   Table suggestion: Leakage audit table (15 rows, one per split)
%   with columns: split, test groups, train count, test count,
%   leaks, status (all should show 0 leaks). Notebook cell 54.


% ------------------------------------------------------------------
\subsection{Preprocessing}\label{subsec:preprocessing}
% ------------------------------------------------------------------
%
% Thatcher R3: open with the claim.
% "Three transformations happen inside the CPCV loop, fitted on
% training data only, to prevent information leakage from test
% fold statistics entering the training pipeline."
%
% This is a per-fold pipeline. Emphasize that each of the 15 splits
% gets its own d*, its own scaler, and its own feature selection.
% The reader must understand that nothing from the test fold leaks
% into these steps.
%
%
% ── Opening paragraph ─────────────────────────────────────────────
%
% Topics:
%   - The preprocessing pipeline runs independently for each of the
%     15 CPCV splits, producing fold-specific transformations
%   - Three sequential steps: fractional differentiation (stationarity),
%     feature scaling (normalization), and feature selection (dimensionality
%     reduction)
%   - Each step is fitted exclusively on the training fold's data;
%     the test fold receives the same transformation using the
%     training-derived parameters (d*, scaler statistics, selected
%     feature list)
%   - This per-fold architecture prevents two forms of leakage:
%     (a) test-fold statistics influencing transformations (e.g.,
%         scaling by test-fold mean/std)
%     (b) feature selection informed by test-fold performance
%
%
% ── 3.6.1 Fractional Differentiation (FFD) ───────────────────────
%
% Topics:
%   - Motivation: before training any model, features should be
%     stationary (no unit root); non-stationary features produce
%     spurious correlations that don't generalize
%   - An ADF test across all 55 features at the 5% significance
%     level identifies ATR as the only non-stationary feature
%     (notebook cell 38); all other features (returns, ratios,
%     bounded oscillators, rate-of-change measures) are already
%     stationary by construction
%   - FFD is therefore applied to ATR only; this is an empirical
%     finding, not an arbitrary design choice
%
%   - FFD procedure (AFML Chapter 5, Snippets 5.1-5.3):
%     (a) Compute FFD weights recursively:
%         w_0 = 1, w_k = -w_{k-1} * (d - k + 1) / k
%         Truncate when |w_k| < tau = 1e-4 or k >= 200
%     (b) Apply as a convolution: FFD(x_t) = sum(w_k * x_{t-k})
%     (c) The first K-1 values become NaN (lookback warm-up);
%         these rows are dropped from both train and test
%
%   - d* estimation (per fold, training data only):
%     (a) Sweep d from 0.0 to 1.0 in steps of 0.05
%     (b) For each d, apply FFD to the training fold's ATR series,
%         then run an ADF test (maxlag=14, autolag="AIC")
%     (c) Select the minimum d where ADF p-value < 0.05
%     (d) If no d achieves stationarity, default to d* = 1.0
%         (full integer differentiation)
%
%   - Application scope: FFD is applied to the FULL ATR series
%     (not just the training fold), using the training-derived d*;
%     this ensures test observations have valid lookback history
%     from preceding (training) values
%   - Anti-leakage argument: FFD is a strictly backward-looking
%     convolution; each output depends only on current and past
%     input values; the only parameter estimated from training
%     is the scalar d*, which does not encode any test-fold
%     information
%
%   - Non-FFD columns: remaining NaN values in non-FFD features
%     (from external data availability gaps or rolling window
%     warm-up) are forward-filled then back-filled within each
%     fold independently
%
%   Cite: \cite{lopez_de_prado_2018}, Chapter 5;
%   \cite{slepaczuk_bieganowski_2024} for walk-forward FFD validation
%
%
% ── 3.6.2 Feature Scaling ────────────────────────────────────────
%
% Topics:
%   - RobustScaler from sklearn: transforms each feature to
%     (x - median) / IQR, where median and IQR are computed from
%     the training fold only
%   - Choice over StandardScaler: BTC features exhibit fat tails
%     and occasional extreme values (flash crashes, halving events);
%     median/IQR is resistant to these outliers, while mean/std
%     is heavily influenced by them
%   - Applied to all 55 features (including the FFD-transformed ATR)
%   - The fitted scaler object is stored in prep_info for potential
%     inverse transformation and for reproducibility documentation
%   - Test fold is transformed using the training-derived median/IQR;
%     test values may fall outside the training range, which
%     RobustScaler handles gracefully (no clipping)
%
%
% ── 3.6.3 Feature Selection: Multi-Model MDA ─────────────────────
%
% Topics:
%   - Purpose: reduce the 55 features to a stable, informative
%     subset that generalizes across model architectures
%   - Method: Mean Decrease Accuracy (MDA), also known as
%     permutation importance (AFML Chapter 8, Section 8.4)
%   - For each feature, permute its values in the test fold and
%     measure the drop in macro-averaged F1 score; larger drop =
%     more important feature
%
%   - Multi-model design (novel relative to standard single-model MDA):
%     (a) MDA computed independently using a Random Forest
%         (500 trees, balanced class weights, captures nonlinear
%         feature interactions and ensemble effects)
%     (b) MDA computed independently using a Logistic Regression
%         (balanced class weights, captures linear effects)
%     (c) Final MDA = average of RF MDA and LR MDA per feature
%     (d) Rationale: a feature ranks high only if it demonstrably
%         contributes to both a linear and a nonlinear classifier;
%         this prevents the selection from being biased toward any
%         single model architecture (e.g., RF might favor
%         interaction-heavy features that LR cannot exploit, and
%         vice versa)
%
%   - Inner CV for MDA: purged 3-fold cross-validation
%     (a) The training fold is divided into 3 chronological blocks
%     (b) For each inner fold, purging is applied using the same
%         t1-based overlap conditions as the outer CPCV
%     (c) This ensures the MDA estimates are not inflated by
%         information leakage within the selection step itself
%
%   - Selection rule (two-stage):
%     Stage 1: keep all features with averaged MDA > 0 (permuting
%       them hurts performance on average across both models)
%     Stage 2: cap at TOP_K_FRAC = 0.40 (top 40% of total features
%       by MDA rank); with 55 features, this caps at 22 features
%     Minimum floor: if fewer than 5 features pass Stage 1,
%       the top 5 by MDA value are kept regardless
%
%   - AR Logistic exception: the AR Logistic model bypasses MDA
%     selection entirely; it constructs its own features (lagged
%     log returns) and receives the pre-selection DataFrame with
%     the log_returns column intact
%
%   - Output: a list of selected feature names, used to subset
%     X_train and X_test before model training; the list is stored
%     in prep_info for stability analysis across folds (Section 3.9)
%
%   Cite: \cite{lopez_de_prado_2018}, Chapter 8 (Section 8.4)
%
%   Table suggestion: Example MDA ranking from one fold showing
%   feature name, MDA_RF, MDA_LR, MDA (average), and selected
%   (yes/no). This illustrates how the dual-model approach works
%   in practice.


% ------------------------------------------------------------------
\subsection{Models}\label{subsec:models}
% ------------------------------------------------------------------
%
% Present all six models grouped by family. A summary table (Table X)
% at the end lists architecture, key hyperparameters, and what each
% model tests. Prose in three subsubsections covers what's unique
% about each family.
%
% Cochrane: "Simple is better. The less math used, the better."
% Only describe what differs across models; shared elements (sample
% weights, class balancing, calibration) are covered in Section 3.8.
%
%
% ── Opening paragraph ─────────────────────────────────────────────
%
% Topics:
%   - Six models are evaluated, spanning four families: econometric,
%     linear ML, ensemble, and neural
%   - This range tests whether model complexity translates into
%     predictive value (RQ2): the simplest model (AR Logistic)
%     uses 6 lagged features and trains in seconds; the most complex
%     (LSTM) uses 21-day sliding windows over 22 selected features
%     and trains in minutes per fold
%   - All models implement a common interface (BaseModel) with
%     methods for fit, predict_proba, predict_logits, and get_name;
%     this ensures identical training/evaluation conditions across
%     architectures
%   - All models receive AFML sample weights and balanced class
%     weights; the integration differs by framework (sklearn
%     sample_weight parameter vs PyTorch per-sample loss multiplier)
%     but the effect is equivalent
%   - Table X summarizes the six models; the subsections below
%     describe what is unique about each family


\subsubsection{Econometric Baseline: AR Logistic}\label{subsubsec:ar}

% Topics:
%   - This model tests a specific hypothesis: can simple price
%     momentum, captured by lagged log returns, match or exceed
%     the predictive power of 55 engineered features?
%   - Feature construction: the model builds its own input from
%     lagged log returns at [1, 2, 3, 5, 10, 21] calendar days,
%     producing 6 features
%   - The lag set covers multiple timescales: daily (1), weekly (2-3),
%     bi-weekly (5), monthly (10), and quarterly-proxy (21)
%   - The longest lag (21) matches the LSTM sliding window length,
%     ensuring both models access the same temporal depth
%   - The model bypasses the MDA-selected feature set entirely;
%     it receives the pre-selection DataFrame and extracts the
%     log_returns column to construct its own lagged features
%   - Architecture: sklearn LogisticRegression with C=1.0 (tunable),
%     L2 penalty (tunable to L1), class_weight="balanced",
%     solver="lbfgs", max_iter=1000
%   - NaN rows from the longest lag (21 days of warm-up) are
%     dropped from training; prediction uses forward/back-fill
%     for edge rows
%   - If this model matches or beats the feature-engineered models,
%     the entire feature engineering pipeline (55 features, 3 external
%     data sources, MDA selection) adds no value beyond simple
%     price momentum for BTC direction prediction
%
%   Cite: Standard econometric baseline; no specific citation needed


\subsubsection{Classical ML Models}\label{subsubsec:classical}

% Topics:
%
% ── Logistic Regression ──────────────────────────────────────────
%
%   - Standard logistic regression on the MDA-selected feature set
%   - The simplest model that uses the full feature engineering pipeline;
%     serves as the linear ML baseline
%   - If ensemble or neural models cannot beat it, nonlinear
%     complexity adds no value for this prediction task
%   - class_weight="balanced": adjusts the loss function to
%     compensate for class imbalance, if present
%   - sample_weight: AFML return-attribution weights (Section 3.3)
%     passed to sklearn's .fit()
%   - Dual weighting: both class_weight and sample_weight are active
%     simultaneously; class_weight adjusts for label imbalance,
%     sample_weight adjusts for label informativeness
%   - Hyperparameters tuned via Optuna (Section 3.8.3):
%     C (regularization strength), penalty (L1 or L2)
%   - Solver: "liblinear" for L1, "lbfgs" for L2 (selected
%     automatically based on the tuned penalty)
%
%
% ── Random Forest ────────────────────────────────────────────────
%
%   - Ensemble of 500 decision trees trained via bagging with
%     random feature subsets at each split
%   - Captures nonlinear feature interactions through ensemble
%     averaging; reduces variance relative to individual trees
%   - class_weight="balanced_subsample": computes class weights
%     per bootstrap sample, adapting to each tree's resampled
%     distribution
%   - max_features="sqrt" (default): each split considers
%     sqrt(n_features) candidates, de-correlating trees
%   - min_samples_leaf=5: prevents deep splits on few observations,
%     providing implicit regularization
%   - Hyperparameters tuned: max_depth (controls tree complexity),
%     min_samples_leaf, max_features; n_estimators capped at 500
%     rather than tuned (more trees rarely hurts, and tuning
%     tree count wastes budget better spent on structural params)
%
%   Cite: \cite{breiman_2001}
%
%
% ── XGBoost ──────────────────────────────────────────────────────
%
%   - Gradient-boosted ensemble of 300 trees trained sequentially;
%     each tree corrects the residual errors of the previous
%   - Objective: binary:logistic (outputs calibrated probabilities)
%   - Early stopping: the calibration set (Section 3.8.1) serves as
%     the evaluation set; training stops after 20 rounds with no
%     improvement in eval_metric (logloss), preventing overfitting
%   - Note: this means the calibration set is "touched" during
%     XGBoost training for round-count selection (but not for
%     feature or weight estimation); this is acknowledged as a
%     mild dependency and is standard practice in boosting
%   - Hyperparameters tuned: max_depth (capped at 6 to prevent
%     memorization), learning_rate, subsample, colsample_bytree,
%     gamma (minimum loss reduction), reg_alpha (L1), reg_lambda (L2)
%   - min_child_weight tuned to control minimum instance weight
%     per leaf, interacting with the sample weights
%
%   Cite: \cite{chen_guestrin_2016}


\subsubsection{Neural Models}\label{subsubsec:neural}

% Topics:
%
% ── LSTM ─────────────────────────────────────────────────────────
%
%   - Architecture: 2-layer LSTM with hidden_size=64, dropout=0.2
%     between layers, followed by a dropout layer and a linear
%     classification head (hidden_size -> n_classes)
%   - Input: sliding window of LSTM_WINDOW=21 days creates 3D
%     sequences of shape (T - 20, 21, n_features); each sequence
%     captures one trading month of feature history
%   - The sliding window reduces the effective sample count: from
%     ~N training observations to ~N-20 sequences; the
%     last_valid_indices array tracks which original indices
%     correspond to each sequence, enabling correct re-alignment
%     of predictions with test timestamps and returns
%   - Optimizer: Adam (lr=1e-3), with ReduceLROnPlateau scheduler
%     (factor=0.5, patience=5); halves the learning rate when
%     validation loss plateaus
%   - Training: mini-batch (batch_size=64), shuffled; up to
%     LSTM_EPOCHS=100 epochs with early stopping (patience=10
%     epochs without validation loss improvement)
%   - Loss: CrossEntropyLoss with class weights (inversely
%     proportional to class frequency) and per-sample AFML weights
%     applied as multipliers: L = mean(w_i * CE(logits_i, y_i))
%   - Best model state (lowest validation loss) is restored after
%     training completes
%   - Hyperparameters tuned: hidden_size (capped at 64), num_layers,
%     dropout, learning_rate
%
%   Cite: \cite{hochreiter_schmidhuber_1997}
%
%
% ── KAN ──────────────────────────────────────────────────────────
%
%   - Architecture: efficient-kan library (Blealtan implementation),
%     which wraps KAN as a standard PyTorch nn.Module
%   - Default architecture: [n_features, 5, 2] (input, one hidden
%     layer with 5 nodes, binary output)
%   - Each edge is parameterized as a B-spline with grid_size=5
%     and spline_order=3 (cubic); the spline is defined over
%     grid_range=[-1, 1]
%   - Input normalization: tanh normalization fitted on training data
%       z = tanh((x - mean_train) / (std_train + 1e-8))
%     This maps features into [-1, 1], matching the spline grid_range;
%     note this is a second normalization on top of RobustScaler
%     (Section 3.6.2), specific to KAN's spline parameterization
%   - Optimizer: AdamW (lr=1e-3, weight_decay=1e-4); the L2
%     penalty prevents spline coefficient explosion
%   - Training: full-batch (all training observations per step);
%     justified by the small effective sample size (~700 training
%     observations per fold after the 80/20 calibration split);
%     up to KAN_EPOCHS=200 with early stopping (patience=20)
%   - Loss: same dual-weighted CrossEntropyLoss as LSTM
%   - Dual-library strategy: efficient-kan is used for all 15 CPCV
%     folds (fast, stable, standard nn.Module); PyKAN is used
%     separately for symbolic extraction (Section 3.10) because
%     only PyKAN exposes the prune(), suggest_symbolic(),
%     fix_symbolic(), and symbolic_formula() APIs
%   - Both libraries share the same B-spline basis and architecture
%     notation; the difference is in training API and symbolic
%     manipulation capabilities
%   - Hyperparameters tuned: width1 (1st hidden layer, range [3, 12]),
%     width2 (2nd hidden layer, range [0, 10]; 0 = single hidden
%     layer), grid size ({3, 5}), learning rate, weight_decay
%
%   Cite: \cite{liu_kan_2024}
%
%
% ── Shared neural design ─────────────────────────────────────────
%
% Topics:
%   - Both LSTM and KAN use dual weighting in the loss function:
%     (a) Class weights: inversely proportional to class frequency,
%         normalized to sum to n_classes; compensates for imbalance
%     (b) Sample weights: AFML return-attribution weights from
%         Section 3.3, applied as per-sample multipliers in the
%         CrossEntropyLoss
%   - Seed count: LSTM and KAN use 2 seeds per fold; AR Logistic,
%     LR, RF, and XGBoost use 3 seeds per fold
%   - Justification for the asymmetry: neural models are
%     computationally more expensive (LSTM: ~36 min total across
%     15 folds; KAN: ~25 min); 2 seeds provide enough variance
%     estimation while keeping total runtime under 2 hours;
%     classical models are cheap enough for 3 seeds
%   - Split-level classification metrics are averaged across
%     available seeds per model; the evaluation code silently
%     handles the asymmetry by skipping missing seed entries
%   - All models use the same reproducibility protocol: fixed
%     random seeds (torch.manual_seed, np.random.seed, random.seed)
%     and deterministic CUDA settings when available
%
%
% ── Summary Table (Table X) ──────────────────────────────────────
%
% Self-contained caption per Cochrane: "each table should have a
% self-contained caption so that a skimming reader can understand
% the fact presented without having to go searching through the text."
%
% Caption: "Summary of the six models evaluated under CPCV. All
% models receive AFML sample weights and balanced class weights.
% Hyperparameters listed under 'Tuned' are optimized per fold via
% Optuna TPE (Section 3.8.3); 'Fixed' parameters are held constant
% across all folds. The Role column describes what hypothesis each
% model tests relative to the research questions."
%
% Columns: Model | Family | Architecture | Fixed Params |
%          Tuned Params | Seeds | Role
% Rows:
%   AR Logistic | Econometric | LR on lagged log_returns
%     [1,2,3,5,10,21] | max_iter=1000 | C, penalty |
%     3 | Tests pure price momentum (RQ2)
%   Logistic Regression | Linear ML | LR on selected features |
%     max_iter=1000 | C, penalty | 3 | Linear baseline (RQ2)
%   Random Forest | Ensemble | 500 trees, balanced_subsample |
%     n_estimators=500 | max_depth, min_leaf, max_features |
%     3 | Nonlinear ensemble (RQ2)
%   XGBoost | Ensemble | 300 trees, early stopping |
%     n_estimators=300, early_stop=20 | max_depth, lr, subsample,
%     colsample, gamma, reg_alpha, reg_lambda | 3 |
%     Gradient boosting (RQ2)
%   LSTM | Neural | 2-layer, window=21, batch=64 |
%     epochs=100, patience=10 | hidden, layers, dropout, lr |
%     2 | Temporal dependencies (RQ2)
%   KAN | Neural | [n, 5, 2], grid=5, k=3, full-batch |
%     epochs=200, patience=20 | width1, width2, grid, lr, wd |
%     2 | Interpretable architecture (RQ1, RQ2)


% ------------------------------------------------------------------
\subsection{Training Pipeline}\label{subsec:training}
% ------------------------------------------------------------------
%
% This section describes the three components of the per-fold
% training procedure that are shared across all models: the
% train/calibration split, probability calibration, and
% hyperparameter tuning. These are presented as a single subsection
% because they form one integrated pipeline, not three independent steps.
%
%
% ── Opening paragraph ─────────────────────────────────────────────
%
% Topics:
%   - For each of the 15 CPCV splits, the training fold undergoes
%     three steps after preprocessing (Section 3.6):
%     (a) An 80/20 chronological split separating model training
%         data from a calibration holdout
%     (b) Nested hyperparameter tuning on the training portion
%     (c) Probability calibration on the calibration holdout
%   - These three steps, combined with the preprocessing pipeline,
%     constitute the full per-fold training procedure; the test
%     fold is never seen until prediction time
%
%
% ── 3.8.1 Train/Calibration Split ────────────────────────────────
%
% Topics:
%   - The preprocessed training fold (after FFD, scaling, and feature
%     selection) is split chronologically at the 80th percentile:
%     the first 80% becomes the model training set (X_model, y_model,
%     w_model), and the last 20% becomes the calibration set
%     (X_cal, y_cal)
%   - Chronological (not random) split preserves temporal ordering,
%     which is critical for time series data
%   - The calibration set serves three purposes:
%     (a) Fitting the probability calibrator (Section 3.8.2)
%     (b) Providing the evaluation set for XGBoost's early stopping
%     (c) Providing the validation set for LSTM/KAN's early stopping
%         and learning rate scheduling
%   - Approximate sizes per fold: with ~700 training observations
%     (after purging/embargo from ~911 total), the model set contains
%     ~560 observations and the calibration set ~140 observations
%   - The calibration set is never used for feature selection (MDA)
%     or FFD estimation (d* sweep); those use the full training fold
%     before the 80/20 split
%   - Design note: using the calibration set for both early stopping
%     and probability calibration means the calibration data is
%     "seen" during training for stopping decisions; this is a mild
%     dependency (the calibration set influences when training stops,
%     not what the model learns), and is standard practice in the
%     literature
%
%   - Hyperparameter tuning (Section 3.8.3) operates on the full
%     training fold (X_tr_sel) with its own inner purged K-fold,
%     not on the 80/20 split; the best hyperparameters are then
%     applied to the model trained on X_model
%
%
% ── 3.8.2 Probability Calibration ────────────────────────────────
%
% Topics:
%   - Motivation: the AFML bet sizing function (Section 3.9.1)
%     converts predicted probabilities into position sizes via an
%     S-curve; if a model outputs P(up)=0.70 but the true accuracy
%     at that confidence level is 0.55, the resulting bet size is
%     too large; calibration ensures that predicted confidence levels
%     correspond to empirical accuracy
%
%   - Two calibration methods, one per model family:
%
%     Platt scaling (sklearn models: AR Logistic, LR, RF, XGBoost):
%       - Fits a logistic regression mapping from raw log-odds to
%         calibrated probabilities on the calibration set
%       - For models that output 2D probability arrays, log-odds
%         are computed as logits[:, 1] - logits[:, 0] before
%         fitting the sigmoid
%       - Effectively unregularized (C=1e10): the sigmoid mapping
%         is purely data-driven with no shrinkage toward the identity
%       - Calibrated probabilities: platt_model.predict_proba(logits)
%
%     Temperature scaling (PyTorch models: LSTM, KAN):
%       - Learns a single scalar T > 0 that minimizes the negative
%         log-likelihood of softmax(logits / T) on the calibration set
%       - T is optimized via bounded scalar minimization (T in [0.1, 10])
%       - T > 1 softens the probability distribution (increases
%         entropy, reduces overconfidence); T < 1 sharpens it
%       - Calibrated probabilities: softmax(logits / T)
%       - LSTM uses a separate fit_from_logits pathway: because
%         the sliding window produces fewer output sequences than
%         input observations, the calibration set indices must be
%         re-aligned using last_valid_indices before fitting T
%
%   Cite: \cite{platt_1999} for Platt scaling,
%   \cite{guo_temperature_2017} for temperature scaling
%
%
% ── 3.8.3 Hyperparameter Tuning ──────────────────────────────────
%
% Topics:
%   - Architecture: nested per-split tuning; inside each of the 15
%     CPCV splits, an Optuna study explores the hyperparameter space
%     using only the training data of that split
%   - This is the correct nested CV architecture per Lopez de Prado
%     (2018, Chapter 7): the outer CPCV provides unbiased test folds,
%     and the inner tuning loop operates entirely within the training
%     fold; the test fold is never seen during tuning
%   - Each of the 15 splits gets its own optimal hyperparameters,
%     simulating how an algorithm would adapt to new data in live
%     trading (the parameters that work best on 2014-2020 data
%     may differ from those that work best on 2016-2022 data)
%
%   - Inner CV for tuning: Purged K-Fold
%     (a) K=3 chronological folds within the training set
%     (b) Embargo = 10 observations between train and validation
%         folds, matching TBL num_days=10; this ensures the inner
%         tuning loop respects the same anti-leakage discipline
%         as the outer CPCV
%     (c) Each Optuna trial trains the model on 2 inner folds and
%         evaluates on the 3rd, rotating through all 3 combinations;
%         the trial score is the mean log-loss across inner folds
%
%   - Optuna configuration:
%     Sampler: TPE (Tree-structured Parzen Estimator); adaptively
%       concentrates trials in promising regions of the search space
%     Pruner: MedianPruner; terminates underperforming trials after
%       evaluating each inner fold, cutting computation on clearly
%       bad configurations
%     Seed: OPTUNA_SEED=42 for reproducibility
%     Verbosity: WARNING level (suppresses per-trial output)
%
%   - Trial budgets:
%     Classical models (LR, RF, XGBoost): 30 trials per split
%     Neural models (LSTM, KAN): 25 trials per split
%     Justification: neural models are 5-10x more expensive per
%       trial; 25 trials provide sufficient exploration while
%       keeping total tuning time manageable (~25 min for KAN
%       across 15 splits)
%
%   - How tuned parameters are applied: the best_params from each
%     split's Optuna study are written to module-level constants
%     (e.g., tree_models.RF_MAX_DEPTH, kan_model.KAN_HIDDEN)
%     before the model training loop for that split; this is
%     equivalent to constructing the model with those parameters
%     directly
%
%   - DSR/PBO validity: because the test fold is never seen during
%     tuning, and because tuning operates within the training fold
%     using its own anti-leakage inner CV, the DSR and PBO computed
%     from test-fold predictions remain valid measures of
%     out-of-sample performance; the n_trials parameter in DSR
%     counts the number of models compared (6), not the number of
%     Optuna trials (25-30 per split)
%
%   Cite: \cite{akiba_optuna_2019} for Optuna,
%   \cite{lopez_de_prado_2018} Chapter 7 for nested CV validity
%
%   Table suggestion: Search spaces for all models (Appendix C).
%   Columns: Model, Parameter, Range/Candidates, Scale.
%   Include a note that grid search is NOT used; TPE is a Bayesian
%   method that does not exhaustively enumerate the space.


% ------------------------------------------------------------------
\subsection{Evaluation Framework}\label{subsec:evaluation}
% ------------------------------------------------------------------
%
% Thatcher R3: open with the claim.
% "The evaluation framework translates raw model predictions into
% financially meaningful performance metrics, corrected for
% selection bias and multiple testing."
%
% This section describes the full chain from calibrated probabilities
% to final model ranking. It is the bridge between the CPCV pipeline
% (Section 3.5-3.8) and the Results chapter (Section 4).
%
%
% ── Opening paragraph ─────────────────────────────────────────────
%
% Topics:
%   - Raw model output: calibrated probabilities P(up), P(down) for
%     each test-fold observation, per model, per split, per seed
%   - These probabilities pass through five evaluation stages:
%     (a) Bet sizing: convert probabilities into position sizes
%     (b) Strategy returns: compute net returns after transaction costs
%     (c) Path stitching: assemble full-span backtests from CPCV paths
%     (d) Performance metrics: Sharpe, drawdown, DSR, PBO
%     (e) Statistical tests: DeLong pairwise AUC significance
%   - Additionally, stability diagnostics assess whether the
%     preprocessing decisions (feature selection, FFD d*) are
%     consistent across folds
%
%
% ── 3.9.1 Bet Sizing (AFML Chapter 10) ──────────────────────────
%
% Topics:
%   - Purpose: convert calibrated probabilities into signed position
%     sizes that reflect both the predicted direction and the model's
%     confidence in that prediction
%   - This is where calibration (Section 3.8.2) matters: the bet
%     size is a direct function of the probability, so miscalibrated
%     probabilities produce systematically wrong position sizes
%
%   - Procedure:
%     (a) Predicted direction: sign = +1 if P(up) > P(down), else -1
%     (b) Confidence: p = max(P(up), P(down)), the probability
%         assigned to the predicted class; p is in [0.5, 1.0]
%     (c) Z-score transform:
%           z = (p - 0.5) / sqrt(p * (1 - p) + 1e-10)
%         This maps the probability space into a standard normal
%         scale; higher confidence -> larger z
%     (d) S-curve mapping (AFML Eq. 10.3):
%           raw_bet = 2 * Phi(z) - 1
%         where Phi is the standard normal CDF; this produces a
%         smooth mapping from probability to bet size in [0, 1]
%     (e) Minimum threshold: if |raw_bet| < 0.10, set bet to 0
%         (no trade); this prevents the model from trading on
%         near-random predictions where P(up) is close to 0.50
%     (f) Apply direction: signed_bet = sign * |raw_bet|
%     (g) Discretization: round to the nearest value in
%         {0.0, 0.25, 0.50, 0.75, 1.0} for realistic position
%         sizing; a portfolio manager cannot take a 0.37-unit position
%
%   - The bet sizing function acts as an implicit abstention mechanism:
%     low-confidence predictions (P close to 0.5) produce zero or
%     very small bets, while high-confidence predictions produce
%     full-sized positions; this is philosophically aligned with the
%     Conservative Predictions approach (Nabar and Shroff, 2023)
%
%   Cite: \cite{lopez_de_prado_2018}, Chapter 10;
%   \cite{nabar_shroff_2023} for the abstention parallel
%
%   Equation suggestion: present the S-curve as a numbered equation
%   in the thesis; it is the only evaluation-stage equation that
%   the reader needs to understand the financial performance metrics
%
%
% ── 3.9.2 Strategy Returns ──────────────────────────────────────
%
% Topics:
%   - For each test-fold observation, the strategy return is computed
%     as the product of the bet size and the realized label return
%     (the return at barrier touch, stored in bins['ret']):
%       gross_return_t = bet_size_t * label_return_t
%
%   - Transaction costs:
%     (a) Cost per unit of turnover: 0.1% (TRANSACTION_COST = 0.001)
%     (b) Turnover at time t: |bet_size_t - bet_size_{t-1}|
%         (absolute change in position size; first trade's turnover
%         is the initial position size)
%     (c) Transaction cost_t = 0.001 * turnover_t
%     (d) Net return_t = gross_return_t - transaction_cost_t
%   - Justification for 0.1% cost: conservative estimate for BTC
%     trading on major exchanges (Binance, Coinbase Pro); maker fees
%     are typically 0.02-0.10%, taker fees 0.04-0.10%; the 0.1%
%     round-trip assumption accounts for both sides of the trade
%     plus potential slippage
%
%   - Annualization: Sharpe and return metrics use ANNUALIZATION_FACTOR
%     = 365 (BTC trades every calendar day, unlike equities at 252)
%   - Risk-free rate: RISK_FREE_RATE = 0.0 (standard for crypto;
%     no universally accepted risk-free rate exists for digital assets)
%
%
% ── 3.9.3 Path Stitching (AFML Chapter 12.4) ─────────────────────
%
% Topics:
%   - Purpose: assemble full-span backtests by combining test-fold
%     predictions from different splits according to the path matrix
%   - For each of the 5 backtest paths, the procedure:
%     (a) Retrieves the path's group-to-split assignments from the
%         path matrix (e.g., path 0 uses split 3 for group 0,
%         split 7 for group 1, etc.)
%     (b) Extracts the calibrated probabilities, true labels, and
%         realized returns from each assigned split's test-fold
%         predictions
%     (c) Concatenates chronologically and sorts by timestamp
%     (d) Applies bet sizing to the concatenated probabilities
%     (e) Computes strategy returns for the full path
%
%   - Each path covers all N=6 groups exactly once, producing a
%     backtest that spans the entire sample period (2014-2026)
%   - Different paths use predictions from different splits for the
%     same time period, so their performance varies; this variation
%     captures the sensitivity of results to the specific train/test
%     partition
%   - Financial metrics are computed per path using seed=0 (the
%     primary seed); classification metrics are averaged across
%     all available seeds
%
%   - Per-path financial metrics computed:
%     Annualized Sharpe ratio:
%       SR = (mean(r) - rf/365) / std(r) * sqrt(365)
%     Cumulative return: prod(1 + r_t) - 1
%     Annualized return: (1 + cum_ret)^(365/n) - 1
%     Maximum drawdown: max peak-to-trough decline in equity curve
%     Time under water: longest consecutive period below previous peak
%     Win rate: fraction of traded observations with positive return
%     Profit factor: sum(positive returns) / |sum(negative returns)|
%     Average bet size: mean(|bet_size|) among traded observations
%     Skewness and kurtosis of strategy returns
%
%   Cite: \cite{lopez_de_prado_2018}, Section 12.4
%
%
% ── 3.9.4 Deflated Sharpe Ratio (AFML Chapter 14) ────────────────
%
% Topics:
%   - Problem: when multiple models are evaluated, the best-performing
%     model's Sharpe ratio is upward-biased due to selection; even
%     under the null hypothesis (all true Sharpes = 0), the maximum
%     observed Sharpe is positive by chance
%   - DSR corrects for this by testing the observed Sharpe against
%     the expected maximum Sharpe under the null
%
%   - Computation:
%     (a) Observed Sharpe: the median Sharpe across 5 paths for each
%         model; the median is used rather than the mean because it
%         is robust to outlier paths
%     (b) Expected maximum Sharpe under null:
%           E[max(SR)] = sqrt(2 * ln(N_trials)) *
%             (1 - gamma / (2 * ln(N_trials))) +
%             gamma / (2 * sqrt(2 * ln(N_trials)))
%         where gamma = 0.5772 (Euler-Mascheroni constant) and
%         N_trials = 6 (number of models compared)
%     (c) Standard error of the Sharpe, corrected for non-normality:
%           SE = sqrt((1 - skew*SR + (kurt-1)/4 * SR^2) / (n-1))
%         where skew and kurt are pooled from path returns
%     (d) DSR = Phi((observed_SR - E[max(SR)]) / SE)
%
%   - Interpretation:
%     DSR > 0.95: the observed Sharpe is statistically significant
%       at the 5% level after correcting for multiple model testing
%     DSR near 0: the observed Sharpe is consistent with random
%       chance; no model demonstrates genuine predictive ability
%
%   - n_trials = 6 (number of models): this is the correct count
%     for model-selection multiplicity; the Optuna hyperparameter
%     trials (25-30 per split) are NOT counted because they occur
%     within the training fold and do not affect the test-fold
%     Sharpe from which DSR is computed
%
%   Cite: \cite{lopez_de_prado_2018}, Chapter 14;
%   \cite{sharpe_1966} for the original Sharpe ratio
%
%   Equation suggestion: present the DSR as a numbered equation;
%   it is the primary evaluation metric (RQ3) and must be
%   interpretable by the reader
%
%
% ── 3.9.5 Probability of Backtest Overfitting (AFML Ch. 11) ──────
%
% Topics:
%   - Problem: even if a model's in-sample Sharpe ratio is high,
%     it may underperform out-of-sample due to overfitting to the
%     specific training/test partition
%   - PBO measures how likely the in-sample best model is to
%     underperform the median model out-of-sample
%
%   - Computation:
%     (a) Input: the Sharpe matrix S of shape (n_models, n_paths),
%         where S[m, p] is the annualized Sharpe of model m on path p
%     (b) For each combination of IS paths and OOS paths (splitting
%         the n_paths = 5 paths into halves, IS_size = 2):
%         - Compute the mean IS Sharpe per model across IS paths
%         - Identify the IS-best model (highest mean IS Sharpe)
%         - Compute the OOS Sharpe for the IS-best model
%         - Compute the OOS median Sharpe across all models
%         - Record whether IS-best underperforms OOS median
%     (c) PBO = fraction of combinations where IS-best underperforms
%
%   - Interpretation:
%     PBO near 0: in-sample model selection is reliable; the
%       IS-best model tends to perform well OOS
%     PBO = 0.5: random; IS selection is no better than chance
%     PBO > 0.5: IS selection is adversarial; the IS winner tends
%       to be the OOS loser (strong overfitting signal)
%
%   - With n_paths = 5 and IS_size = 2, there are C(5, 2) = 10
%     IS/OOS combinations; PBO resolution is therefore 0.10
%
%   Cite: \cite{lopez_de_prado_2018}, Chapter 11
%
%
% ── 3.9.6 Statistical Significance: DeLong Test ──────────────────
%
% Topics:
%   - Purpose: test whether two models' AUC-ROC values are
%     statistically different on the same test observations
%   - The DeLong test (DeLong et al., 1988) is a non-parametric
%     method for comparing correlated ROC curves
%
%   - Procedure:
%     (a) Pool predictions across all 15 CPCV splits for each model
%         (non-overlapping test sets make pooling valid)
%     (b) For each pair of models (A, B):
%         - Compute AUC_A and AUC_B on the pooled test set
%         - Compute the variance of each AUC and their covariance
%           using DeLong's structural components (placement values)
%         - Z-statistic: z = (AUC_A - AUC_B) / sqrt(var_diff)
%         - Two-sided p-value: p = 2 * Phi(-|z|)
%     (c) Significant at alpha = 0.05 if p < 0.05
%
%   - Total pairwise tests: C(6, 2) = 15 model pairs
%   - Note on LSTM alignment: LSTM produces fewer predictions than
%     other models (due to the sliding window dropping the first 20
%     observations); when comparing LSTM against another model, the
%     shorter array length is used
%   - No multiple testing correction (e.g., Bonferroni) is applied
%     to the 15 pairwise tests; this is common in the CPCV literature
%     but should be acknowledged as a limitation
%
%   Cite: \cite{delong_1988}
%
%   Table suggestion: DeLong test results as a table with columns
%   model_a, model_b, AUC_a, AUC_b, delta_AUC, z_stat, p_value,
%   significant (yes/no)
%
%
% ── 3.9.7 Stability Diagnostics ──────────────────────────────────
%
% Topics:
%   - Two stability metrics assess whether the preprocessing
%     decisions are consistent across the 15 CPCV folds:
%
%   Feature selection stability:
%     - For each feature, count how many of the 15 folds selected
%       it (using seed=0 and the first model's prep_info)
%     - Compute selection frequency = count / 15
%     - Features with frequency > 0.80 (selected in 12+ of 15 folds)
%       are classified as "stable"
%     - Interpretation: if most features have high stability, the
%       MDA selection is identifying a consistent signal; if stability
%       is flat (many features near 50-60%), no compact feature subset
%       dominates, suggesting diffuse or weak signal
%     - A flat stability profile is itself a finding: it suggests
%       that the information relevant to BTC direction prediction is
%       distributed across many features rather than concentrated
%       in a few
%
%   FFD stability:
%     - For each FFD-transformed column (ATR in this thesis), compute
%       the standard deviation of d* across all 15 folds
%     - std(d*) < 0.1: the stationarity structure is consistent
%       across time periods; the same degree of differentiation
%       is needed regardless of which data partition is used
%     - std(d*) > 0.1: the optimal differentiation order changes
%       across folds, indicating that the feature's stationarity
%       properties are time-varying (e.g., ATR might be more
%       persistent during low-volatility regimes and less during
%       high-volatility regimes)
%
%   - Both diagnostics are reported in the Results chapter
%     (Section 4.4) and serve as quality checks on the pipeline's
%     internal consistency
%
%
% ── 3.9.8 Model Comparison and Ranking ───────────────────────────
%
% Topics:
%   - The final model comparison aggregates all metrics into a
%     single ranked table:
%     Primary ranking criterion: median Sharpe ratio (descending)
%     Tiebreaker: standard deviation of Sharpe (ascending; lower
%       variability across paths is preferred)
%   - Additional columns: DSR, mean F1-macro, mean accuracy,
%     mean AUC-ROC, mean Brier score, median max drawdown,
%     median cumulative return, median win rate, median profit factor
%   - The table provides a comprehensive view: classification metrics
%     (F1, AUC) and financial metrics (Sharpe, drawdown) do not
%     always agree; a model can have decent F1 but poor Sharpe if
%     it mispredicts on high-return events
%
%   Table suggestion: Model comparison table (Table X in Results).
%   Self-contained caption describing all metrics and the ranking
%   criterion. This is the "main result" table per Cochrane.


% ------------------------------------------------------------------
\subsection{Symbolic Extraction}\label{subsec:symbolic}
% ------------------------------------------------------------------
%
% Framed as exploratory analysis (per advisor guidance), not the
% core contribution. Keep concise in the main text (~1.5 pages);
% detailed training logs, edge-by-edge R^2 values, and unsimplified
% formulas go to Appendix F.
%
% Thatcher R3: open with the purpose, not the procedure.
% "Symbolic extraction produces a closed-form mathematical decision
% function from a trained KAN, enabling direct inspection of which
% features drive the prediction and how they combine."
%
%
% ── Opening paragraph ─────────────────────────────────────────────
%
% Topics:
%   - The symbolic extraction pipeline operates independently of the
%     CPCV evaluation loop; it is applied after all 15 splits have
%     been evaluated
%   - Purpose: transform the trained KAN's learned B-spline
%     activations into interpretable symbolic functions (e.g., sin,
%     exp, x^2, tanh), producing a closed-form decision function
%     that a human can read, audit, and validate against domain
%     knowledge
%   - This addresses RQ1: can a trained KAN produce interpretable
%     symbolic decision functions for BTC daily direction prediction?
%   - The CPCV pipeline uses efficient-kan (standard nn.Module,
%     fast, stable); symbolic extraction uses PyKAN independently
%     because only PyKAN exposes the prune(), suggest_symbolic(),
%     fix_symbolic(), and symbolic_formula() APIs
%   - Both libraries share the same B-spline basis and architecture
%     notation ([n_features, hidden, n_classes], grid size, spline
%     order k=3); the difference is in training API and symbolic
%     manipulation capabilities
%
%
% ── Fold selection ────────────────────────────────────────────────
%
% Topics:
%   - Symbolic extraction is performed on a single CPCV fold, not
%     on all 15 (PyKAN training is expensive and the extraction is
%     exploratory)
%   - fold_selection="last" (split 14): the fold whose test set
%     covers the most recent data; closest to how a practitioner
%     would deploy the model in a rolling-window scenario
%   - Alternative: fold_selection="best" would select the fold with
%     the highest KAN F1-macro; this maximizes the chance of
%     meaningful symbolic output but biases toward the fold that
%     happened to work best
%   - Justification for "last": the most recent data reflects the
%     current market regime; a formula extracted from 2024-2026 data
%     is more actionable than one extracted from 2018-2020 data
%   - This choice is acknowledged as a limitation: the formula is
%     fold-specific and may change if a different fold is selected
%
%
% ── Data preparation ──────────────────────────────────────────────
%
% Topics:
%   - The selected fold's preprocessing is reconstructed: FFD
%     (using the stored d*), scaling (using the stored RobustScaler),
%     and feature selection (using the stored feature list)
%   - An optional feature subset can further restrict the features
%     to the top N most stable (by CPCV selection frequency); this
%     thesis uses n_top_features=5 to produce a simpler, more
%     interpretable formula
%   - The training fold is split 80/20 (same as the CPCV pipeline)
%     to provide a validation set for early stopping
%   - Tanh normalization is applied: z = tanh((x - mean) / (std + 1e-8)),
%     fitted on the training portion, applied to both; this matches
%     both efficient-kan and PyKAN's expected input range [-1, 1]
%
%
% ── PyKAN training (Algorithm 1, Step 1) ──────────────────────────
%
% Topics (cite: cho_lee_kim_2025 for Algorithm 1):
%   - A fresh PyKAN model is trained from scratch on the selected
%     fold's data (not a transfer from efficient-kan)
%
%   Architecture selection (priority order):
%     1. If tuned CPCV params are available for the selected fold
%        (from the Optuna study in Section 3.8.3), the tuned width1
%        is used as the starting point, with three simplifications
%        for symbolic tractability:
%        (a) Width cap: width1 is capped at PYKAN_SYMBOLIC_WIDTH_CAP=8;
%            a width-12 KAN produces ~24 summed terms pre-pruning,
%            which overwhelms sympy during formula simplification
%        (b) Drop width2: the second hidden layer is removed
%            (PYKAN_SYMBOLIC_DROP_WIDTH2=True); two hidden layers
%            produce nested f(g(x)) compositions that sympy cannot
%            simplify in reasonable time
%        (c) Force coarse grid: grid is set to PYKAN_SYMBOLIC_FORCE_GRID=3
%            regardless of the tuned grid; coarser splines produce
%            cleaner symbolic matches with fewer basis functions
%     2. If no tuned params exist for this fold, falls back to
%        data-aware sizing: hidden width chosen so that
%        n_train / total_params >= 5 (PYKAN_MIN_SAMPLES_PER_PARAM)
%     In both cases, the final width is additionally clamped by the
%     data-aware safety floor to prevent memorization; with ~358
%     training observations and ~6 parameters per edge (grid + k),
%     this typically yields a hidden width of 3-5
%
%   - This design keeps the symbolic model architecturally linked
%     to the benchmarked CPCV KAN (same width1 origin) while
%     ensuring formulas remain tractable for sympy
%
%   - Three-phase training:
%     Phase 1 (Adam, 600 steps): main learning phase; weight decay
%       = 1e-3; Gaussian noise (std=0.05) injected into inputs each
%       step and clipped to [-1, 1]; the noise acts as a dropout-like
%       regularizer for small datasets; best validation state saved
%     Phase 2a (LBFGS warmup, ~20 steps): light refinement with no
%       sparsity regularization; best validation state saved
%     Phase 2b (LBFGS sparsity, ~20 steps): L1 + entropy
%       regularization (lamb=0.002) encourages sparse activations,
%       which aids downstream pruning; early stopping patience=10
%   - Grid extension is disabled: with ~358 training samples, refining
%     the grid from 3 to 5 would add parameters and cause memorization;
%     grid stays at 3 throughout
%   - Minimum accuracy gate: if validation accuracy after Adam is
%     below 53% (barely above random for binary classification),
%     a warning is logged; the pipeline continues but symbolic
%     output is likely to yield constant functions
%
%
% ── Pruning (Algorithm 1, Step 2) ─────────────────────────────────
%
% Topics:
%   - After training, edges with activation magnitude below
%     PRUNE_THRESHOLD=0.01 are removed
%   - The model.attribute() method computes per-edge importance scores
%   - model.prune() removes dead edges and collapses nodes that lost
%     all connections
%   - Diagnostic: pre-prune and post-prune edge counts and survival
%     rates are logged; very few active edges (< 3) indicate that
%     regularization was too aggressive
%   - Post-prune accuracy is checked: if it drops significantly,
%     pruning removed signal-carrying edges
%   - The pruned architecture (e.g., [5, 3, 2] → [5, 2, 2]) is
%     recorded and reported in the results
%
%
% ── Symbolification (Algorithm 1, Step 3) ─────────────────────────
%
% Topics:
%   - For each surviving edge (l, i, j), the B-spline activation
%     is replaced with a closed-form symbolic function
%   - Candidate library (14 functions):
%     x, x^2, x^3, x^4, exp, log, sqrt, tanh, sin, cos, abs,
%     sgn, arctan, 0 (constant)
%   - For each edge, model.suggest_symbolic() returns candidate
%     functions ranked by a combination of R^2 fit and complexity
%   - Selection criterion: accept the best non-constant candidate
%     with R^2 >= 0.30 (SYMBOLIC_R2_THRESHOLD)
%   - Brute-force fallback: when suggest_symbolic() returns "0"
%     (constant function wins by total_loss due to zero complexity),
%     the pipeline iterates through all 14 candidates manually,
%     calls fix_symbolic() for each, parses the R^2 from PyKAN's
%     printed output, and selects the best non-constant candidate
%     above the R^2 threshold; model state is restored between
%     each candidate trial
%   - R^2 threshold of 0.30 (lowered from 0.50): empirical choice;
%     0.50 was too aggressive for the weak-signal regime and rejected
%     most edges, producing near-constant formulas; 0.30 allows more
%     edges to be symbolified while maintaining a minimum fit quality
%   - Symbolification rate: the fraction of total edges that are
%     successfully replaced with symbolic functions; reported in
%     the results alongside the edge-by-edge R^2 distribution
%
%
% ── Affine fine-tuning (Algorithm 1, Step 4) ──────────────────────
%
% Topics:
%   - After symbolification, the symbolic function's affine
%     parameters (scale and shift: a*f(b*x + c) + d) are fine-tuned
%     via 30 LBFGS steps at lr=0.0004 (from the VIX KAN paper)
%   - This adjusts the symbolic fit without changing the function
%     family (e.g., "tanh" stays "tanh" but its scale/shift may change)
%   - If fine-tuning produces NaN loss, the pre-fine-tune state is
%     restored
%
%
% ── Formula extraction ────────────────────────────────────────────
%
% Topics:
%   - model.symbolic_formula() returns sympy expressions for each
%     output node
%   - For binary classification, two expressions are returned:
%     logit_bearish (class 0) and logit_bullish (class 1)
%   - The decision function is: decision = logit_bullish - logit_bearish
%   - The probability of an up move is:
%     P(up) = 1 / (1 + exp(-decision)) = sigmoid(decision)
%   - Variable substitution: PyKAN uses internal variable names
%     (x_1, x_2, ...); these are replaced with the actual feature
%     names (e.g., macd_hist, yield_curve_10y30y) via sympy.subs()
%   - Variable naming convention: the pipeline tries 1-based (x1),
%     1-based with underscore (x_1), and 0-based (x_0), selecting
%     whichever convention PyKAN used internally
%   - sympy.simplify with a 30-second timeout: complex expressions
%     can cause sympy to hang; if the timeout triggers, the
%     unsimplified expression is used
%   - sympy.nsimplify with tolerance=1e-3 cleans up floating-point
%     coefficients into rational approximations
%
%   - Output dictionary:
%     decision_function: the sympy expression (human-readable formula)
%     p_up_formula: "1 / (1 + exp(-(decision)))"
%     surviving_features: list of feature names that appear as free
%       symbols in the decision function (features that were pruned
%       or symbolified to constants disappear from the formula)
%     pre_symbolic_accuracy: validation accuracy before symbolification
%     post_symbolic_accuracy: validation accuracy after symbolification
%     symbolification_rate: fraction of edges symbolified
%     pruned_architecture: the architecture after pruning
%
%
% ── Limitations (acknowledged here, discussed in Section 5.3) ────
%
% Topics:
%   - Fold-specific: the formula depends on which fold is selected;
%     different folds may produce different formulas with different
%     surviving features
%   - Small sample: ~358 training observations after the 80/20 split;
%     this limits the PyKAN architecture (hidden width constrained
%     by the 5 samples/parameter rule) and prevents grid extension
%   - Architectural simplifications: the symbolic extraction KAN is
%     not identical to the CPCV-benchmarked KAN; width1 is capped
%     at 8 (vs tuned values up to 12), width2 is dropped (single
%     hidden layer forced), and grid is forced to 3 (vs tuned 5);
%     these simplifications are necessary for sympy tractability
%     but mean the symbolic formula approximates a simplified
%     version of the CPCV model, not the exact model that produced
%     the benchmarked predictions
%   - R^2 threshold sensitivity: lowering the threshold from 0.50
%     to 0.30 increases the symbolification rate but accepts
%     lower-quality symbolic fits
%   - Post-symbolic accuracy may be lower than pre-symbolic accuracy
%     because the symbolic functions approximate but do not perfectly
%     reproduce the learned B-spline activations
%   - The brute-force fallback introduces computational cost and
%     depends on PyKAN's internal output format, which varies
%     across versions
%   - These limitations are consistent with the advisor's guidance
%     to frame symbolic extraction as exploratory rather than as
%     the core contribution
%
%   Cite: \cite{cho_lee_kim_2025} for Algorithm 1,
%   \cite{liu_kan_2024} for the symbolic extraction concept,
%   \cite{practitioner_kan_2025} for training best practices
%
%   Appendix F: full PyKAN training logs, phase-by-phase accuracy,
%   edge-by-edge R^2 values with function assignments, pruned
%   architecture diagrams, unsimplified formulas before nsimplify


% ==================================================================
% CHAPTER 4: RESULTS
% ==================================================================
\newpage
\section{Results}\label{sec:results}

% Cochrane: "Start with the main result."
% Thatcher R3: "Lead with the contribution."
% Thatcher R10: "Frame empirics as theory tests."
%
% Structure: main result table first (4.1), then supporting evidence
% in decreasing order of importance. Each subsection answers a
% specific RQ or provides evidence for the Discussion chapter.
%
% Placeholders use [VALUE] for numbers to be filled in from the
% actual pipeline run.
%
%
% ── Opening paragraphs (before any subsection) ───────────────────
%
% Topics:
%   - State the headline finding immediately:
%     "No model achieves a Deflated Sharpe Ratio above [VALUE],
%      indicating that none of the six architectures demonstrates
%      statistically significant predictive ability for BTC daily
%      direction after correcting for multiple model testing."
%   - Follow with the secondary finding:
%     "The symbolic extraction pipeline successfully produces a
%      closed-form decision function with [VALUE]% symbolification
%      rate, identifying [FEATURE_1], [FEATURE_2], and [FEATURE_3]
%      as the features that survive in the final formula."
%   - Brief roadmap: Section 4.1 presents the model comparison,
%     4.2 the classification detail, 4.3 the financial performance,
%     4.4 the stability diagnostics, and 4.5 the symbolic results


\subsection{Model Comparison}\label{subsec:results_comparison}

% This is the "main result" per Cochrane: the table that answers
% RQ2 (does KAN justify its complexity?) and RQ3 (does AFML change
% conclusions?). Present it first, discuss it briefly, save
% interpretation for Chapter 5.
%
%
% ── Main comparison table (Table X) ───────────────────────────────
%
% Topics:
%   - Present the full model comparison table described in Section 3.9.8
%   - Ranked by median Sharpe ratio (descending)
%
%   Self-contained caption:
%   "Model comparison across six architectures evaluated under
%   CPCV (N=6, k=2, 15 splits, 5 backtest paths). Median Sharpe
%   is computed across 5 paths. DSR corrects for multiple model
%   testing (n_trials=6). Classification metrics (F1-macro, AUC,
%   accuracy) are averaged across 15 splits and all available seeds.
%   All models use AFML sample weights, balanced class weights,
%   and per-split Optuna hyperparameter tuning."
%
%   Columns and placeholder values:
%     Model | Med. Sharpe | Sharpe Std | DSR | F1-macro | AUC |
%     Accuracy | Brier | Rank
%
%     AR Logistic      | [VALUE] | [VALUE] | [VALUE] | [VALUE] | [VALUE] | [VALUE] | [VALUE] | [RANK]
%     Log. Regression  | [VALUE] | [VALUE] | [VALUE] | [VALUE] | [VALUE] | [VALUE] | [VALUE] | [RANK]
%     Random Forest    | [VALUE] | [VALUE] | [VALUE] | [VALUE] | [VALUE] | [VALUE] | [VALUE] | [RANK]
%     XGBoost          | [VALUE] | [VALUE] | [VALUE] | [VALUE] | [VALUE] | [VALUE] | [VALUE] | [RANK]
%     LSTM             | [VALUE] | [VALUE] | [VALUE] | [VALUE] | [VALUE] | [VALUE] | [VALUE] | [RANK]
%     KAN              | [VALUE] | [VALUE] | [VALUE] | [VALUE] | [VALUE] | [VALUE] | [VALUE] | [RANK]
%
%
% ── Key observations (state facts, save interpretation for Ch. 5) ─
%
% Topics:
%   - DSR result: all models show DSR = [VALUE], below the 0.95
%     significance threshold; no model demonstrates statistically
%     significant predictive performance after correcting for
%     the number of models tested
%   - Ranking: [BEST_MODEL] achieves the highest median Sharpe
%     ([VALUE]), followed by [2ND] ([VALUE]) and [3RD] ([VALUE]);
%     KAN ranks [KAN_RANK]th with median Sharpe [VALUE]
%   - Classification vs financial metrics: note any divergence
%     (e.g., a model with higher AUC but lower Sharpe indicates
%     that its errors concentrate on high-return events)
%   - The AR Logistic baseline [matches/underperforms/outperforms]
%     the feature-engineered models, suggesting that the 55
%     engineered features [do/do not] add value beyond simple
%     price momentum
%   - All metrics should use 2-3 significant digits per Cochrane


\subsection{Classification Metrics}\label{subsec:results_classification}

% Supporting detail for the model comparison. This section provides
% the evidence behind the summary metrics in Table X.
%
%
% ── Per-split classification summary ──────────────────────────────
%
% Topics:
%   - For each model, report the distribution of F1-macro across
%     15 splits (mean, std, min, max)
%   - Note whether performance is consistent across splits or
%     highly variable; high variability suggests regime dependence
%     (the model works in some time periods but not others)
%
%   Table suggestion: Per-split F1 distribution table
%   Columns: Model | Mean F1 | Std F1 | Min F1 | Max F1
%     AR Logistic      | [VALUE] | [VALUE] | [VALUE] | [VALUE]
%     Log. Regression  | [VALUE] | [VALUE] | [VALUE] | [VALUE]
%     Random Forest    | [VALUE] | [VALUE] | [VALUE] | [VALUE]
%     XGBoost          | [VALUE] | [VALUE] | [VALUE] | [VALUE]
%     LSTM             | [VALUE] | [VALUE] | [VALUE] | [VALUE]
%     KAN              | [VALUE] | [VALUE] | [VALUE] | [VALUE]
%
%
% ── AUC-ROC ──────────────────────────────────────────────────────
%
% Topics:
%   - Pooled AUC per model (pooled across all 15 splits; valid
%     because test sets are non-overlapping)
%   - Report pooled AUC values:
%     AR Logistic: [VALUE], LR: [VALUE], RF: [VALUE],
%     XGBoost: [VALUE], LSTM: [VALUE], KAN: [VALUE]
%   - Note whether any model's AUC significantly exceeds 0.50
%     (random classifier baseline)
%
%
% ── DeLong pairwise AUC test ──────────────────────────────────────
%
% Topics:
%   - Report the DeLong test results for the [NUMBER] significant
%     pairs out of 15 total pairwise comparisons
%   - If few or no pairs are significant, state this explicitly:
%     "No pairwise AUC difference reaches significance at
%     alpha = 0.05, indicating that all models perform similarly
%     in discriminative ability"
%   - If some pairs are significant, report which models differ
%     and the direction of the difference
%
%   Table suggestion: DeLong results (only significant pairs, or
%   full 15-pair table if space permits)
%   Columns: Model A | Model B | AUC_A | AUC_B | z | p-value | Sig.
%
%
% ── Confusion matrices (optional, can go to appendix) ─────────────
%
% Topics:
%   - Aggregated confusion matrix per model (summed across all 15
%     splits, seed 0): TP, FP, TN, FN
%   - Derive precision, recall per class from the confusion matrix
%   - Note whether any model exhibits a strong bias toward one class
%     (e.g., predicting "Up" for nearly all observations)
%   - If all models show near-random class distributions (close to
%     50/50 predicted), this is consistent with the DSR=0 finding
%   - Detailed per-split confusion matrices go to Appendix D
%
%   Table suggestion: Aggregated confusion matrix per model
%   or a compact table with columns:
%   Model | Precision(Up) | Recall(Up) | Precision(Down) | Recall(Down)


\subsection{Financial Performance}\label{subsec:results_financial}

% This section presents the financial metrics that answer the
% practical question: "Would a trader make money using these models?"
% The answer (spoiler: no) is the evidence for RQ3.
%
%
% ── Path-level Sharpe distribution ────────────────────────────────
%
% Topics:
%   - For each model, report the Sharpe ratio across the 5 backtest
%     paths: median, mean, std, min, max
%   - Highlight whether any path achieves a positive Sharpe; if so,
%     note that the DSR still rejects significance because the
%     positive paths are offset by negative ones (or because the
%     maximum is insufficient after n_trials correction)
%
%   Table suggestion: Path-level Sharpe table
%   Columns: Model | Path 0 | Path 1 | Path 2 | Path 3 | Path 4 |
%     Median | Std
%     AR Logistic      | [V] | [V] | [V] | [V] | [V] | [V] | [V]
%     Log. Regression  | [V] | [V] | [V] | [V] | [V] | [V] | [V]
%     Random Forest    | [V] | [V] | [V] | [V] | [V] | [V] | [V]
%     XGBoost          | [V] | [V] | [V] | [V] | [V] | [V] | [V]
%     LSTM             | [V] | [V] | [V] | [V] | [V] | [V] | [V]
%     KAN              | [V] | [V] | [V] | [V] | [V] | [V] | [V]
%
%
% ── DSR computation detail ────────────────────────────────────────
%
% Topics:
%   - Report the DSR inputs for the best-performing model:
%     observed median Sharpe = [VALUE]
%     expected max Sharpe under null (n_trials=6) = [VALUE]
%     pooled skewness = [VALUE], pooled kurtosis = [VALUE]
%     SE(SR) = [VALUE]
%     DSR = [VALUE]
%   - Repeat for all models (or present as a single table)
%   - Key statement: "DSR = [VALUE] for all models, falling well
%     below the 0.95 threshold for statistical significance"
%
%
% ── PBO result ────────────────────────────────────────────────────
%
% Topics:
%   - Report PBO = [VALUE] (out of 10 IS/OOS combinations)
%   - Interpretation:
%     If PBO near 0: in-sample model selection is reliable, but
%       the selected model still has DSR=0 (reliable selection of
%       a model that doesn't work is not useful)
%     If PBO near 0.5: random; model selection adds no value
%     If PBO > 0.5: overfitting; the IS winner underperforms OOS
%   - Note that PBO near 0 combined with DSR=0 is the most
%     intellectually honest outcome: it means the pipeline is not
%     overfitting (good methodology), but also not finding signal
%     (the market is efficient relative to these features)
%
%
% ── Equity curves ─────────────────────────────────────────────────
%
% Topics:
%   - Plot cumulative strategy return for each model across the
%     best path (or median path)
%   - Overlay BTC buy-and-hold cumulative return for comparison
%   - Note whether any model's equity curve stays above 1.0
%     (positive cumulative return) for sustained periods, or
%     whether all curves oscillate around the zero-return line
%
%   Figure suggestion: Equity curves for all 6 models + buy-and-hold,
%   one panel per model or all on the same chart (depends on
%   readability). Notebook cell [NUMBER].
%
%
% ── Additional financial metrics (compact table) ──────────────────
%
% Topics:
%   - For each model (median across paths):
%     Cumulative return: [VALUE]%
%     Annualized return: [VALUE]%
%     Max drawdown: [VALUE]%
%     Time under water: [VALUE] days
%     Win rate: [VALUE]%
%     Profit factor: [VALUE]
%     Average bet size: [VALUE]
%     Fraction of observations traded (bet != 0): [VALUE]%
%
%   Table suggestion: Financial performance summary
%   Self-contained caption: "Financial performance metrics for each
%   model, computed as the median across 5 CPCV backtest paths.
%   Strategy returns include 0.1% per-trade transaction costs and
%   are annualized at 365 days."


\subsection{Feature Selection Stability}\label{subsec:results_features}

% This section reports the MDA feature selection consistency across
% the 15 CPCV folds and the FFD d* stability. It does not directly
% answer an RQ but provides evidence for the Discussion chapter
% (Section 5.1: interpreting negative results).
%
%
% ── Feature selection frequency ───────────────────────────────────
%
% Topics:
%   - Report the selection frequency (count / 15) for each feature
%   - State the headline finding:
%     If flat: "No feature is selected in all 15 folds. The most
%       frequently selected feature is [FEATURE] ([COUNT]/15 = [FREQ]%).
%       [NUMBER] features are selected in fewer than 5 folds. The
%       distribution is relatively flat, suggesting that no compact
%       feature subset consistently dominates across time periods."
%     If concentrated: "A stable core of [NUMBER] features is selected
%       in 12 or more of 15 folds: [FEATURE_1], [FEATURE_2], ...
%       This suggests these features carry consistent signal."
%   - Group the findings by feature type:
%     TA features: [NUMBER] selected in >10 folds, [NUMBER] in <5
%     Math features: [NUMBER] selected in >10 folds, [NUMBER] in <5
%     External features: [NUMBER] selected in >10 folds, [NUMBER] in <5
%     On-chain features specifically: [COMMENT on whether any
%       on-chain feature shows high stability]
%   - Connect to RQ1: the features that survive symbolic extraction
%     (Section 4.5) should ideally overlap with the most frequently
%     selected features here; if they don't, this is worth discussing
%
%   Figure suggestion: Horizontal bar chart of selection frequency
%   (55 bars, sorted by frequency, colored by feature group:
%   blue=TA, green=math, orange=external). Notebook cell [NUMBER].
%
%   Table suggestion: Top 10 and bottom 10 features by selection
%   frequency (compact table for the main text; full 55-feature
%   table in Appendix E)
%   Columns: Rank | Feature | Group | Frequency (count/15) | %
%
%
% ── FFD stability ─────────────────────────────────────────────────
%
% Topics:
%   - Report the d* values across the 15 folds for ATR:
%     Mean d* = [VALUE], Std d* = [VALUE], Min = [VALUE], Max = [VALUE]
%   - Interpretation:
%     If std < 0.1: "The optimal differentiation order for ATR is
%       stable across folds (std = [VALUE]), indicating consistent
%       stationarity structure regardless of the training period."
%     If std > 0.1: "The differentiation order varies across folds
%       (std = [VALUE]), suggesting that ATR's persistence properties
%       change across market regimes (e.g., more persistent during
%       low-volatility periods, less during high-volatility episodes)."
%   - Report the most common d* value (mode) and whether any folds
%     required d* = 1.0 (full differentiation, meaning ATR is
%     heavily persistent in that period)
%
%   Table suggestion: d* per fold (15 rows)
%   Columns: Split | Test Groups | d* | ADF p-value at d*
%   Or a compact summary: Mean | Std | Min | Max | Mode


\subsection{Symbolic Extraction Results}\label{subsec:results_symbolic}

% This section answers RQ1. Framed as exploratory analysis per
% advisor guidance. Present what the pipeline produced, then
% acknowledge limitations; interpretation goes to Section 5.3.
%
%
% ── Extraction summary ────────────────────────────────────────────
%
% Topics:
%   - Fold used: split [NUMBER] (fold_selection="last")
%   - Test group date range: [START_DATE] to [END_DATE]
%   - Number of MDA-selected features entering extraction: [NUMBER]
%   - After n_top_features=5 restriction: [FEATURE_1], [FEATURE_2],
%     [FEATURE_3], [FEATURE_4], [FEATURE_5]
%   - Training sample size (after 80/20 split): [NUMBER] observations
%   - Data-aware architecture: [INPUT, HIDDEN, OUTPUT]
%     (e.g., [5, 3, 2])
%
%
% ── Training diagnostics ─────────────────────────────────────────
%
% Topics:
%   - Phase 1 (Adam, 600 steps):
%     Final train loss = [VALUE], val loss = [VALUE]
%     Val accuracy = [VALUE]%
%   - Phase 2a (LBFGS warmup):
%     Val accuracy = [VALUE]%
%   - Phase 2b (LBFGS sparsity):
%     Val accuracy = [VALUE]%
%   - State whether the minimum accuracy gate (53%) was passed
%   - Pre-prune accuracy: [VALUE]%
%
%   Note: detailed phase-by-phase training logs go to Appendix F
%
%
% ── Pruning results ───────────────────────────────────────────────
%
% Topics:
%   - Pre-prune edge count: [NUMBER]
%   - Post-prune edge count: [NUMBER]
%   - Survival rate: [VALUE]%
%   - Pruned architecture: [INPUT, HIDDEN_PRUNED, OUTPUT]
%     (e.g., [5, 3, 2] → [5, 2, 2])
%   - Post-prune accuracy: [VALUE]% (note any drop from pre-prune)
%
%
% ── Symbolification results ───────────────────────────────────────
%
% Topics:
%   - Total edges attempted: [NUMBER]
%   - Successfully symbolified: [NUMBER]
%   - Symbolification rate: [VALUE]%
%   - Edge-by-edge results:
%     Edge (layer, in, out) | Suggested function | R^2 | Method
%     (l, i, j)            | [FUNCTION]          | [V] | suggest/brute
%     ...
%   - If any edges used the brute-force fallback, note how many
%     and which functions were selected
%   - If any edges fell below R^2 = 0.30, note they were left as
%     constant ("0") and effectively removed from the formula
%
%   Table suggestion: Edge symbolification table
%   Columns: Edge | Function | R^2 | Method (suggest/brute-force)
%
%
% ── The decision function ─────────────────────────────────────────
%
% Topics:
%   - Present the extracted formula:
%     P(up) = sigmoid([FORMULA])
%     where [FORMULA] is the sympy-simplified decision function
%   - Example format (placeholder):
%     P(up) = sigmoid(a * tanh(b * [FEATURE_1] + c) +
%              d * sin(e * [FEATURE_2] + f) + g)
%     where a, b, c, d, e, f, g are the numerical coefficients
%   - State clearly which features appear as free symbols in the
%     decision function: the "surviving features"
%     Surviving features: [FEATURE_1], [FEATURE_2], [FEATURE_3]
%   - Note any features that entered the extraction but were pruned
%     or symbolified to constants (effectively removed)
%
%   - Present the formula as a numbered equation in the thesis;
%     this is the concrete output of RQ1 and should be immediately
%     visible to a skimming reader
%
%
% ── Pre-symbolic vs post-symbolic accuracy ────────────────────────
%
% Topics:
%   - Pre-symbolic accuracy (trained KAN before symbolification):
%     [VALUE]%
%   - Post-symbolic accuracy (after replacing splines with symbolic
%     functions and affine fine-tuning): [VALUE]%
%   - Delta: [VALUE] percentage points
%   - Interpretation:
%     If delta is small (< 3 pp): "The symbolic functions approximate
%       the learned splines closely; the formula preserves most of
%       the KAN's predictive structure"
%     If delta is large (> 5 pp): "The symbolic approximation degrades
%       accuracy by [VALUE] pp, indicating that some learned B-spline
%       shapes are not well captured by the 14-function library; the
%       formula should be interpreted as an approximation, not an
%       exact reproduction of the KAN's decision boundary"
%   - Compare post-symbolic accuracy to random baseline (50%):
%     If above 55%: the formula retains meaningful signal
%     If near 50%: the formula is effectively random despite being
%       well-structured (the underlying signal is too weak for
%       symbolic distillation to capture)
%
%
% ── Feature interpretation (brief, details in Section 5.3) ────────
%
% Topics:
%   - For each surviving feature, note:
%     (a) Which feature group it belongs to (TA, math, external,
%         on-chain)
%     (b) Its selection frequency across 15 folds (from Section 4.4)
%     (c) The functional form the KAN assigns to it (e.g., tanh,
%         sin, x^2)
%   - Note whether on-chain features appear in the formula (relevant
%     to Problem 1 from the Introduction: does the "public inside
%     view" of blockchain data provide measurable signal?)
%   - State the observation; save the economic interpretation for
%     Discussion (Section 5.3):
%     "[FEATURE_1] enters the formula via [FUNCTION], suggesting
%      a [monotonic/periodic/nonlinear] relationship with BTC
%      direction; [FEATURE_2] enters via [FUNCTION], suggesting..."
%   - If no on-chain features survive: note this explicitly as a
%     finding (the on-chain "inside view" does not provide enough
%     signal for the KAN to retain it after pruning)
%
%
% ── Comparison with best-fold extraction (optional) ───────────────
%
% Topics:
%   - If the pipeline was also run with fold_selection="best"
%     (split [NUMBER]), briefly compare:
%     Best fold symbolification rate: [VALUE]%
%     Best fold surviving features: [FEATURE_LIST]
%     Best fold post-symbolic accuracy: [VALUE]%
%   - Note differences: do the same features survive? Does the
%     functional form change? Is the symbolification rate higher?
%   - If the formulas differ substantially, this illustrates the
%     fold-specificity limitation acknowledged in Section 3.10
%   - If this comparison was not run, skip this block
%
%
% ── Section closing ───────────────────────────────────────────────
%
% Topics:
%   - Summarize the RQ1 answer factually:
%     "The symbolic extraction pipeline produces a closed-form
%      decision function with [VALUE]% symbolification rate. The
%      formula identifies [NUMBER] surviving features: [LIST].
%      Post-symbolic accuracy is [VALUE]%, [above/near] the random
%      baseline. The procedure's limitations, including fold
%      specificity and sample size constraints, are discussed in
%      Section 5.3."


% ==================================================================
% CHAPTER 5: DISCUSSION
% ==================================================================
\newpage
\section{Discussion}\label{sec:discussion}

% Interpret findings. Link back to mechanisms (Thatcher R7).

\subsection{Interpreting the Negative Results}\label{subsec:disc_negative}
% DSR=0 as consistent with EMH and market efficiency.
% HARd to Beat: properly fitted baselines are hard to surpass.
% Conservative Predictions: abstention beats prediction in low-signal regimes.
% Three-version pipeline evolution as methodological maturity narrative.

\subsection{KAN Performance in Context}\label{subsec:disc_kan}
% KAN fourth, behind RF and XGBoost but ahead of LSTM.
% Contrast with KASPER's strong results (different asset, regime detection).

\subsection{Symbolic Extraction as Contribution}\label{subsec:disc_symbolic}
% Interpretability value even when predictive performance is weak.
% Comparison with VIX KAN paper's symbolic results.
% Limitations: fold-specific, small sample, R² threshold sensitivity.

\subsection{Limitations}\label{subsec:disc_limitations}
% Sample size (911 observations), daily frequency,
% single asset, seed asymmetry, fold_selection="last" vs "best".


% ==================================================================
% CHAPTER 6: CONCLUSION
% ==================================================================
\newpage
\section{Conclusion}\label{sec:conclusion}

% Target: max 10% of textual part (~3.5 pages).
% Short. State the contribution. End with resonance (Thatcher R16).
% Do NOT restate all findings (Cochrane).


% ==================================================================
% CHAPTER 7: FUTURE WORK
% ==================================================================
\newpage
\section{Future Work}\label{sec:future}

% Keep brief (Cochrane: don't write your grant application here).
% Concrete, actionable directions only.
%
% Candidates:
% - MultKAN (KAN 2.0) comparison for multiplicative interactions
% - Higher-frequency data (hourly/4h bars)
% - Alternative assets (ETH, gold) for symbolic formula comparison
% - Expanding window vs rolling window CPCV
% - Meta-labeling layer on top of best classifier
% - Larger sample periods or multi-asset pooling


% =============================================
% BIBLIOGRAPHY
% =============================================
\newpage
\phantomsection
\thispagestyle{plain}
\addcontentsline{toc}{section}{Bibliography}
\bibliographystyle{agsm}
\bibliography{references}
% Create a references.bib file with all your citations.


% =============================================
% APPENDICES
% =============================================
\newpage
\phantomsection
\thispagestyle{plain}

\appendix
\section{Appendices}\label{sec:appendices}

\subsection{Pipeline Architecture Diagram}\label{app:pipeline}
% Full pipeline flowchart: data → labeling → weights → features
% → alignment → CPCV → preprocessing → models → evaluation → symbolic.

\subsection{CPCV Split Details}\label{app:cpcv_splits}
% Group boundaries, train/test timelines, purging/embargo verification.

\subsection{Hyperparameter Search Spaces}\label{app:hyperparams}
% Optuna search spaces for all models (table format).

\subsection{Per-Split Classification Reports}\label{app:split_reports}
% Detailed per-split metrics if needed.

\subsection{Full Feature List}\label{app:features}
% Complete table of 55 features with descriptions, sources, parameters.

\subsection{Symbolic Extraction Detailed Output}\label{app:symbolic}
% Full PyKAN training logs, edge-by-edge R² values,
% pruned architecture diagrams, unsimplified formulas.


% =============================================
% End Document
% =============================================
\end{document}